\documentclass{article}
\usepackage{graphicx} % Required for inserting images
\usepackage{amsmath}
\usepackage{amsfonts}
\usepackage{booktabs}
\usepackage[hidelinks]{hyperref}

\title{cs312\_a2}
\author{suzeva }
\date{October 2026}

\begin{document}

\maketitle

\section{Problem 1: Learning-rate scaling}
\subsection{(a) Learn from small training budgets}

I used the nine supplied d8 runs: peak learning rate
$\eta\in\{0.0015,0.003,0.006\}$ at
$D\in\{153.6,307.2,614.4\}$M training tokens.  All nine runs use the same
model, so this part scales the training budget $D$, not model size.

At each budget I fit a quadratic in log learning rate to approximate the
loss minimum (Figure~\ref{fig:p1a}, left).
Define
\[
x=\log_2(\eta/0.003),
\qquad
L_D(\eta)=a_Dx^2+b_Dx+c_D.
\]
The sampled learning rates are $x\in\{-1,0,1\}$.  This fits the U-shaped
loss curve; a simple power law would be monotone.  The fitted minimum is
$x^*=-b_D/(2a_D)$ and $\eta^*=0.003\,2^{x^*}$.

\begin{figure}[ht]
    \centering
    \includegraphics[width=\textwidth]{figures/p1a_lr_fits_and_scaling.png}
    \caption{Left: measured final losses (circles), quadratic fits in log LR,
    and their minima (stars).  Right: a power-law fit to the three inferred
    optimal learning rates. }
    \label{fig:p1a}
\end{figure}

\begin{table}[ht]
\centering
\small
\begin{tabular}{rrrrrr}
\toprule
$D$ (tokens) & $a_D$ & $b_D$ & $c_D$ & fitted $\eta^*$ & fitted loss \\
\midrule
153.6M & 0.045784 &  0.030720 & 3.229134 & 0.002378 & 3.223981 \\
307.2M & 0.014959 &  0.000310 & 3.056223 & 0.002979 & 3.056222 \\
614.4M & 0.018412 & $-0.003297$ & 2.925636 & 0.003192 & 2.925488 \\
\bottomrule
\end{tabular}
\caption{The three loss--LR laws and their inferred optima, where
$x=\log_2(\eta/0.003)$.}
\label{tab:p1a-lr-fits}
\end{table}

I then fit a power law relating the training-token budget $D$ to the optimal
learning rate (Figure~\ref{fig:p1a}, right):
\[
\eta^*(D)=\eta_{\mathrm{ref}}
\left(\frac{D}{D_{\mathrm{ref}}}\right)^\beta
=0.003276\left(\frac{D}{614.4\mathrm{M}}\right)^{0.2125}.
\]
Longer training prefers a larger peak LR.  With only three measurements per
budget, precise minima depend on the fit; quadratics in raw LR give a similar
scaling exponent (0.221).  Part (b) tests the prediction.

\paragraph{Takeaway.} Our fitted rule: double the training data, multiply
peak LR by $1.16$.  This estimate uses only three budgets.

\subsection{(b) Predict before increasing the budget}

Before inspecting the provided losses at the larger budgets, I applied the
power law from part (a).  I then fit the provided loss--LR curves using
the same quadratic in log LR (Figure~\ref{fig:p1b}, left):

\begin{center}
\begin{tabular}{lrrr}
\toprule
Training tokens & Predicted LR & Fitted optimum & Error (\%) \\
\midrule
1.2288B & 0.003796 & 0.003878 & $-2.1$ \\
1.8432B & 0.004138 & 0.003426 & $+20.8$ \\
2.4576B & 0.004399 & 0.002906 & $+51.4$ \\
\bottomrule
\end{tabular}
\end{center}

Error is $(\eta_{\mathrm{pred}}/\eta_{\mathrm{fit}}-1)\times100\%$;
positive means the prediction is too high.

\begin{figure}[ht]
    \centering
    \includegraphics[width=\textwidth]{figures/p1b_predictions_and_fits.png}
    \caption{Left: all six budgets from (a) and (b), with provided losses
    (circles), quadratic fits, and fitted minima (stars).
    Right: the recorded predictions (hollow diamonds)
    compared with fitted optima from (a) and (b).  The largest token
    budget is yellow.}
    \label{fig:p1b}
\end{figure}

Figure~\ref{fig:p1b} (right) shows that the first prediction is close,
but the next two overshoot the fitted optima.
The best measured LR is 0.003 at all three budgets.
The fitted optima decrease across these larger budgets; this reversal also
appears with quadratics in raw LR, although their precise minima differ.
The optima between sampled LRs are estimates, not additional measurements.

\paragraph{Takeaway.} The small-budget upward trend does not continue, it reverses:
extrapolating it overshoots the larger-budget optima.

\subsection{(c) Compare two scaling rules}

I fit a straight line in log tokens versus log optimal LR, using either
all six source budgets or only 1.2288B, 1.8432B, and 2.4576B.  The resulting
power laws are
\begin{align*}
\eta_{\mathrm{all}}(D)
&=0.00348169\left(\frac{D}{2.4576\mathrm{B}}\right)^{0.09881},\\
\eta_{\mathrm{large}}(D)
&=0.00295666\left(\frac{D}{2.4576\mathrm{B}}\right)^{-0.40928}.
\end{align*}
I recorded both predictions at 4.9152B tokens before inspecting target losses
(Figure~\ref{fig:p1c-predictions}).

\begin{figure}[ht]
    \centering
    \includegraphics[width=\textwidth]{figures/p1c_scaling_predictions.png}
    \caption{Power laws fitted to all six source optima (left) or only
    the three larger-budget optima (right).  Stars are source optima;
    dashed lines extrapolate to the recorded target LRs (yellow hollow diamonds).}
    \label{fig:p1c-predictions}
\end{figure}

Before seeing target losses, I expected the all-six fit to do better because
more source budgets could reduce variance.  The completed runs instead favor
the larger-three fit:

\begin{center}
\small
\begin{tabular}{lrrr}
\toprule
Run & Peak LR & Final val loss & Gap to grid best \\
\midrule
Grid & 0.00150000 & 2.717728 & $+0.000249$ \\
Grid (best) & 0.00300000 & 2.717479 & $0$ \\
Grid & 0.00600000 & 2.728922 & $+0.011443$ \\
All-six prediction & 0.00372850 & 2.719499 & $+0.002020$ \\
Larger-three prediction & 0.00222637 & 2.716902 & $-0.000576$ \\
\bottomrule
\end{tabular}
\end{center}

The reference is the best sampled LR in the prescribed three-point grid,
not the best of all five runs; negative gaps indicate improvement.
The same quadratic in log LR, fitted only to the three target-grid losses,
estimates $\eta^*=0.002153$ (Figure~\ref{fig:p1c-results}, left).
The larger-three prediction is 3.4\% above
this estimate, versus 73.2\% for the all-six prediction.
Fitting all five measurements gives a similar optimum, $0.002165$.

\begin{figure}[ht]
    \centering
    \includegraphics[width=\textwidth]{figures/p1c_target_results.png}
    \caption{Left: final validation losses at all five LRs (circles),
    the quadratic fitted to the three fixed-grid runs, and its minimum
    (star).  Dashed vertical lines mark the two predicted LRs.
    Right: their measured loss gaps relative to the best fixed-grid run.
    Prediction colors match across panels; fixed-grid points and their
    fit are teal.}
    \label{fig:p1c-results}
\end{figure}

The larger-three rule follows the recent downturn and beats the all-six
rule by 0.002597 loss.  It also improves slightly on the grid best, although
the 0.000576 gain is a small, single-seed difference.

This illustrates the bias--variance trade-off.  Using only three nearby
budgets may increase variance because fewer points determine the rule,
but avoids fitting the earlier trend which might be more biased.  Adding the three smaller
budgets could reduce variance, yet makes a single power law a worse proxy
for the target regime (increasing bias): its extrapolation is biased toward too large an LR.
We cannot measure variance separately from these single-seed runs.

Both rules extrapolate $2\times$ beyond their largest source budget
(2.4576B to 4.9152B).  The larger-three rule gives a small measured gain
at this distance; the all-six rule does not.  Relative to the smallest
source budget, the target is $4\times$ larger for the three-budget fit
and $32\times$ larger for the six-budget fit.  Those are full
source-to-target spans, not different extrapolation distances.
We have not tested how much farther either rule can extrapolate.
All these comparisons change training tokens, not model size.

\paragraph{Takeaway.} We found no gains from including the smaller-budget
fits.  Fewer nearby budgets were a better proxy for the target: less
trend mismatch mattered more than having more points.

\subsection{(d) Hyperball versus AdamW}

I used all 24 supplied Hyperball runs: eight LRs at each of
$D\in\{153.6,\allowbreak307.2,\allowbreak614.4\}$M tokens. At each budget I least-squares fit
the same quadratic in log LR as in (a), using all eight measurements:
$L_D(\eta)=a_Dx^2+b_Dx+c_D$, with $x=\log_2(\eta/.003)$.
Its minimum gives the following estimates, compared with AdamW at the
same budgets:

\begin{center}
\begin{tabular}{lrr}
\toprule
Training tokens & Hyperball fitted LR & AdamW fitted LR \\
\midrule
153.6M & 0.015610 & 0.002378 \\
307.2M & 0.012213 & 0.002979 \\
614.4M & 0.010416 & 0.003192 \\
\bottomrule
\end{tabular}
\end{center}

Fitting a power law to the three Hyperball optima gives
\[
\eta^*_{\mathrm{Hyperball}}(D)
=0.010267\left(\frac{D}{614.4\mathrm{M}}\right)^{-0.2919}.
\]
Figure~\ref{fig:p1d-source} shows Hyperball's fitted optimal LR decreasing
with training budget, opposite to
AdamW's small-budget exponent $+0.2125$ from (a).  Absolute LRs are not
equivalent step sizes because the optimizers scale updates differently.

\begin{figure}[ht]
    \centering
    \includegraphics[width=\textwidth]{figures/p1d_hyperball_and_adamw.png}
    \caption{Left: all supplied Hyperball losses and quadratic fits in
    log LR.  Middle: the same fits zoomed near their minima (stars).
    Right: power laws fitted to Hyperball (filled stars) and AdamW
    (hollow stars) optima at the same three budgets.  The largest token
    budget is yellow.}
    \label{fig:p1d-source}
\end{figure}

The exponent is sensitive to the loss fit: using only the four LRs nearest
the minima ($.006,.01,.015,.03$) gives $-0.1335$, still decreasing.
Only three Hyperball source budgets are supplied; these fits alone
do not establish reliable extrapolation.

Before testing the target, I recorded the prediction at 1.2288B tokens:
\[
\eta^*_{\mathrm{Hyperball}}(1.2288\mathrm{B})
=0.010267\,2^{-0.2919}\approx\mathbf{0.00838685}.
\]

The target runs used our local AdamH implementation, with the same parameter
routing and equivalent update equations as the subsequently supplied course code.
All four runs completed with the following final losses:
\begin{center}
\begin{tabular}{lrr}
\toprule
Peak LR & Final val loss & Gap to best sampled \\
\midrule
0.006000 & 2.843828 & $+0.018086$ \\
0.008387 (prediction) & 2.829280 & $+0.003539$ \\
0.010000 & 2.825742 & $0$ \\
0.012000 & 2.828076 & $+0.002334$ \\
\bottomrule
\end{tabular}
\end{center}

Fitting the three nearby controls, with the prediction run held out,
gives $\eta^*=0.009991$ (Figure~\ref{fig:p1d-target}, left).
The prediction is $16.1\%$ below that optimum;
fitting all four runs instead gives $0.010250$, so the optimum remains
near $0.010$. The target LR still decreases, but less than the source
law predicts.

\begin{figure}[ht]
    \centering
    \includegraphics[width=\textwidth]{figures/p1d_target_results.png}
    \caption{Left: target losses (circles), a quadratic fitted to the
    three nearby controls, and its optimum (star); purple marks the
    prediction run. Right: source-only power law and its extrapolation,
    with the target prediction (hollow diamond) and fitted optimum
    (yellow star).}
    \label{fig:p1d-target}
\end{figure}

Hyperball's three source optima fit a power law more closely than AdamW's
(log-space $R^2=0.985$ versus $0.914$). But its target LR error is larger:
$16.1\%$ versus AdamW's $2.1\%$ at the same budget in (b).
A tight three-point source fit therefore does not establish more
predictable extrapolation. Both target optima are interpolated from
sparse, single-seed measurements.

\paragraph{Takeaway.} Hyperball's optimal LR decreases as the training
budget grows, but its power law did not predict the larger-budget optimum
more accurately than AdamW's. This comparison uses only one target budget.

\subsection{(e) Optional: transfer across LR schedules}

I compare the nine supplied cosine-decay runs with the nine linear-decay
runs from (a), keeping model size, WD, warmup and token budgets fixed.
Both schedules choose $.003$ as the best \emph{sampled} LR at all three
source budgets. Using the same quadratic in $x=\log_2(\eta/.003)$ gives:
\begin{center}
\begin{tabular}{lrr}
\toprule
Training tokens & Linear fitted LR & Cosine fitted LR \\
\midrule
153.6M & .002378 & .002256 \\
307.2M & .002979 & .002534 \\
614.4M & .003192 & .003191 \\
\bottomrule
\end{tabular}
\end{center}
Figure~\ref{fig:p1e-source} shows the measured curves and optima.
The source-only cosine rule is
\[
\eta^*_{\mathrm{cos}}(D)=.0031305
\left(\frac{D}{614.4\mathrm{M}}\right)^{.2500},
\]
versus exponent $.2125$ for linear decay. Doubling the tokens therefore
multiplies LR by about $1.19$ for cosine, versus $1.16$ for linear.
This rule predicts $.0037228$ at 1.2288B tokens, close to the linear
prediction $.0037961$. I recorded it before submitting four new cosine
target runs at $\{.0015,.003,.0037228,.006\}$. The best point was at the
upper endpoint, so I added $.012$ to bracket it; all five runs completed.
No supplied source runs were repeated.

\begin{figure}[htbp]
\centering
\includegraphics[width=\linewidth]{figures/p1e_schedule_source.png}
\caption{Optional schedule comparison. Left/middle: supplied loss--LR
measurements (circles), quadratic fits and interpolated optima (stars).
The largest source budget is yellow. Right: source-only power laws,
solid within the fitted range and dashed beyond it, with frozen target
predictions (hollow diamonds) and held-out fitted optima (yellow stars).}
\label{fig:p1e-source}
\end{figure}

Figure~\ref{fig:p1e-target} compares the completed target curves. At
1.2288B tokens, the local three-point cosine fit gives optimal LR
$.004798$, versus $.003878$ for the supplied linear curve. The cosine
prediction has loss $2.837759$, only $.000430$ above the best sampled
cosine loss $2.837329$ at LR $.006$. Directly transferring LR $.003$
gives $2.840015$, a gap of $.002686$.
The precise cosine optimum is uncertain: fitting all five points instead
gives $.004059$. This changes the estimated best LR more than it changes
the practical conclusion that the frozen prediction is near the loss minimum.

\begin{figure}[htbp]
\centering
\includegraphics[width=\linewidth]{figures/p1e_schedule_target.png}
\caption{Schedule transfer at 1.2288B tokens. Measured losses are circles;
stars mark local quadratic optima. Hollow outlined circles identify direct
LR $.003$; the cosine prediction is a hollow diamond. The cosine fit uses
the best point and its two neighbors, not all five measurements.}
\label{fig:p1e-target}
\end{figure}

The best sampled cosine and linear losses differ by only $-.000460$;
these single-seed, unequally dense grids do not establish a schedule winner.
At the same LR $.003$, cosine is $.002226$ worse than linear. Retuning
can therefore change the ranking seen at a single fixed LR.

\paragraph{Takeaway.} Linear and cosine decay use similar LR scales.
The cosine scaling prediction transfers well in loss, despite uncertainty
in the exact optimum. Neither schedule has a clear tuned-loss advantage here.

\section{Problem 2: Joint learning-rate and weight-decay scaling}
\subsection{(a) Let LR and WD vary together}

I used the 36 supplied runs: nine LR--WD combinations at each of four
training-token budgets, with model size fixed.  For each budget I least-squares fit
\[
\widehat L(x,y)=a+bx+cy+dx^2+exy+fy^2,
\qquad
x=\log_2(\eta/0.003),\quad y=\log_2(\lambda/0.1).
\]
Here $\lambda$ is weight decay; the $exy$ term captures LR--WD interaction.
All four fitted surfaces have a minimum inside the sampled grid
(Figure~\ref{fig:p2a-contours}).

\begin{figure}[ht]
    \centering
    \includegraphics[width=\textwidth]{figures/p2a_joint_contours.png}
    \caption{Quadratic loss contours for each budget, with measured losses
    (circles) and fitted optima (stars).  Both axes are logarithmic.
    Color denotes loss, with a separate scale for each panel.
    Dashed white lines hold $\eta\lambda=\eta^*\lambda^*$ fixed (part b).}
    \label{fig:p2a-contours}
\end{figure}

\begin{center}
\begin{tabular}{lrr}
\toprule
Training tokens & Fitted LR & Fitted WD \\
\midrule
153.6M & 0.001859 & 0.8466 \\
307.2M & 0.001603 & 0.5194 \\
614.4M & 0.002470 & 0.2579 \\
1.2288B & 0.002696 & 0.1622 \\
\bottomrule
\end{tabular}
\end{center}

\begin{figure}[ht]
    \centering
    \includegraphics[width=\textwidth]{figures/p2a_optima_and_loss_gains.png}
    \caption{Left: joint-fit optimal LR and Problem 1's fixed-WD LR optima.
    Middle: fitted optimal WD. Both include power-law fits.
    Right: best measured losses from both grids; labels show
    $\Delta L=L_{\mathrm{P1,best}}-L_{\mathrm{joint,best}}$.
    Stars denote fitted optima; circles denote measured losses.
    The largest token budget is yellow.}
    \label{fig:p2a-trends}
\end{figure}

\paragraph{Takeaways.}
\begin{enumerate}
\item Figure~\ref{fig:p2a-trends} shows that tuning WD gives a lower
optimal peak LR than fixing WD at $.1$.
LR still rises overall with training tokens, though it dips at 307.2M.
\item Optimal WD falls with training tokens. The fitted exponent is
$-0.816$: doubling the tokens multiplies optimal WD by $2^{-0.816}\approx0.568$.
\item A single power law describes WD much better than LR:
log-space $R^2=0.994$ versus $0.678$.
\item WD tuning helps most at small budgets. The measured loss gain
shrinks from $0.0447$ to $0.0017$ (about $0.002$) at 1.2288B tokens.
That last gain is small; without repeated seeds, we cannot establish
a reliable benefit or estimate run-to-run variance.
\end{enumerate}

\subsection{(b) From coupling to a joint scaling law}

On each contour in Figure~\ref{fig:p2a-contours}, I drew
$\lambda=\eta^*\lambda^*/\eta$: increasing LR decreases WD along this
line. It roughly follows the low-loss trade-off, most clearly at
1.2288B tokens. Compensation keeps loss near its minimum only over a
limited range; farther along the line, loss increases even though the
product stays fixed.

Using the four products of the fitted optima from part (a), I fit the
power law shown in Figure~\ref{fig:p2b-product}:
\[
\eta^*(D)\lambda^*(D)
=0.00041963\left(\frac{D}{1.2288\mathrm{B}}\right)^{-0.59294}.
\]
Doubling the training tokens multiplies the fitted product by $0.663$.
All three power-law comparisons use least squares in log space and
the same four source budgets.

\begin{figure}[ht]
    \centering
    \includegraphics[width=.8\textwidth]{figures/p2b_product_scaling.png}
    \caption{Products of fitted optima (stars) and their power-law fit.
    The inset compares log-space $R^2$ for LR, WD, and their product.
    The largest token budget is yellow.}
    \label{fig:p2b-product}
\end{figure}

\paragraph{Takeaway.} WD alone still has the tightest power-law scaling,
but LR$\times$WD follows a power law much better than LR alone.
This comparison uses only four single-seed source budgets.

\subsection{(c) Test the joint rule}

Before inspecting target results, the product law from (b) predicted
\[
\lambda_{\mathrm{pred}}
=\frac{(\eta^*\lambda^*)(2.4576\mathrm{B})}{0.003}
=0.092736.
\]
The 2.4576B-token budget was excluded from the fit. I tested this WD
at peak LR $.003$ and compared it with the best measured 153.6M pair
transferred unchanged, and a WD sweep at the target. All runs use AdamW
and linear decay. Four runs are new; the $.003$ LR, $.1$ WD control
is supplied by P1(b).

\begin{center}
\begin{tabular}{lrrr}
\toprule
Recipe & Peak LR & WD & Final val loss \\
\midrule
Product-rule prediction & 0.003 & 0.092736 & 2.768667 \\
Transferred small-budget pair & 0.0015 & 1.6 & 2.852842 \\
WD-grid control & 0.003 & 0.05 & 2.771656 \\
WD-grid control (supplied) & 0.003 & 0.1 & 2.768251 \\
WD-grid control & 0.003 & 0.2 & 2.773303 \\
\bottomrule
\end{tabular}
\end{center}

The product prediction lowers loss by $.084174$ relative to the
transferred small-budget pair. Its loss is $.000416$ higher than the
best WD-grid control.
The latter difference is small; these single-seed results do not establish
a reliable advantage for either nearby WD.

\paragraph{Takeaway.} The product rule predicted an effective recipe
without target-budget tuning. The baseline WD sweep independently found
a very similar recipe ($\lambda=.1$); the rule nearly matched it, but did
not improve on its measured loss.

\section{Problem 3.1: Batch size in the noisy quadratic model}

All experiments run locally on CPU with $N=8{,}192$ examples and exactly
$N/B$ updates. The initial coordinates have variances $1$ and $.1$,
as specified in the handout. I use constant LR, no WD, and standard Adam with
$\beta_2=.95$, $\epsilon=10^{-8}$; RMSProp is Adam with $\beta_1=0$.
Source LR sweeps average 2,048 independent trajectories on a coarse grid
and 4,096 on a refined grid; batches $B\geq128$ use 8,192 and 32,768.
A local quadratic in log LR through three
neighboring points interpolates the minimum. Final losses use 32,768 fresh
trajectories, independent of tuning; prediction and tuned-LR comparisons
share initialization/noise to reduce comparison variance.

\subsection{(a) SGD}

Fitting the seven estimated optima at $B\in\{1,2,4,8,16,32,64\}$
(Figure~\ref{fig:p31a}, middle) gives
\[
\eta^*(B)=0.036874\left(\frac{B}{64}\right)^{0.99206}.
\]
Doubling the batch multiplies LR by $1.989$, almost exactly two.
I froze this rule before sweeping $B=256$: it predicts $0.145882$,
versus a tuned optimum of $0.123795$. Their final losses are
$0.000700$ and $0.000552$, respectively.

\begin{figure}[ht]
    \centering
    \includegraphics[width=\textwidth]{figures/p31a_sgd.png}
    \caption{Left: selected SGD sweeps (circles), exact expected-loss curves
    (lines), and interpolated optima (stars). Middle: the source power law
    and its extrapolation; the hollow diamond is the batch-256 prediction.
    Right: independently evaluated loss after LR tuning, with approximate
    95\% confidence intervals. Larger batches are yellow.}
    \label{fig:p31a}
\end{figure}

The rule remains useful through $B=128$ (about $1.5\%$ excess expected
loss), but at 256 its LR is $18\%$ too high and its loss is $27\%$ higher.
At 512 it predicts $0.2902$, beyond SGD's stability limit
$\eta<2/\lambda_{\max}(H)=0.2$.

Doubling $B$ halves the gradient-noise variance, but also halves the
number of updates at fixed $N$. For small steps, doubling LR compensates:
the total optimization progress scales with $\eta N/B$, while the noise
floor scales with $\eta/B$, so both stay roughly unchanged. This explains
the near-linear LR rule and initially flat tuned loss. At larger batches,
raising LR amplifies noise in the high-curvature direction and approaches
the stability limit. Fewer stable updates also leave more initialization
error, so tuned loss rises despite less noisy gradients.

\paragraph{Takeaway.} Double the batch, roughly double LR---but only while
steps remain small. Larger batches give cleaner gradients but fewer chances
to update; eventually LR cannot safely increase enough to compensate.

\subsection{(b) RMSProp and Adam}

The \emph{source exponent} is the fitted power $p$ in $\eta^*(B)=cB^p$,
using only the seven source batch sizes $B=1$--$64$, before testing at 256.
It measures how quickly optimal LR increases with batch size:
doubling the batch multiplies LR by $2^p$.
The left panels of Figure~\ref{fig:p31b} show the RMSProp and Adam
loss--LR sweeps. Fitting their source optima gives
\[
\begin{aligned}
\eta^*_{\mathrm{RMSProp}}(B)&=0.025831(B/64)^{0.90423},\\
\eta^*_{\mathrm{Adam}}(B)&=0.025456(B/64)^{0.90497}.
\end{aligned}
\]
Their exponents are nearly identical ($p\approx0.904$--$0.905$), slightly
below SGD's $0.992$: both adaptive optimizers suggest multiplying LR by
$1.87$ when doubling the batch, versus SGD's $1.99$. Despite this agreement,
their frozen predictions transfer differently to $B=256$:

\begin{center}
\small
\begin{tabular}{lrrrr}
\toprule
Optimizer & Predicted LR & Tuned LR & Predicted loss & Tuned loss \\
\midrule
RMSProp & 0.090478 & 0.093144 & 0.013001 & 0.012946 \\
Adam & 0.089256 & 0.175288 & 0.015175 & 0.013954 \\
\bottomrule
\end{tabular}
\end{center}

RMSProp's loss penalty is about $0.4\%$ (gap $0.000055$, paired
SE $0.000038$); Adam's is $8.8\%$ (gap $0.001222$, paired
SE $0.000120$). Adam's target loss curve is relatively flat and irregular,
so its estimated optimal LR is less precise than its loss comparison.

\begin{figure}[ht]
    \centering
    \includegraphics[width=\textwidth]{figures/p31b_adaptive.png}
    \caption{Left: cached RMSProp and Adam loss--LR sweeps near their
    minima; lines join measured points. Batch colors increase to yellow
    at 256. Middle: solid power laws fit only batches 1--64; dashed lines
    extrapolate beyond 64, and hollow diamonds mark the batch-256
    predictions. Stars mark interpolated optima.
    Right: independently evaluated tuned losses (top) and predicted-versus-tuned
    losses at 256 (bottom, approximate 95\% confidence intervals).
    SGD is included in the middle/right comparisons; its sweeps appear
    in Figure~\ref{fig:p31a}.}
    \label{fig:p31b}
\end{figure}

\paragraph{Takeaway.} RMSProp and Adam follow nearly the same scaling rule
over the source batches, with almost identical exponents. This does not
imply similar transfer accuracy at larger batches.

\paragraph{Language-model prediction.} At fixed training tokens, I expect
optimal LR to increase with batch size, approximately following a power
law over $B=8$--$64$. I expect extrapolation to 128 to be more reliable
than to 256, with LR growth eventually slowing as fewer updates and
stability become limiting. These are hypotheses to test in Problem 3.2,
not established batch-size cutoffs; I would fit the LM exponent rather
than assume the toy model's $p\approx0.905$ transfers unchanged.

\subsection{(c) Noise and curvature}

\paragraph{Starting setting.}
Part (b) uses $f(w_1,w_2)=(w_1^2+10w_2^2)/2$, with independent
$w_{0,1}\sim\mathcal N(0,1)$ and $w_{0,2}\sim\mathcal N(0,.1)$.
The noisy gradient is $g_t=Hw_t+\epsilon_t$, where
$H=\operatorname{diag}(1,10)$ and $\epsilon_t\sim\mathcal N(0,I_2/B)$.
Training uses $N=8{,}192$ examples, $N/B$ updates, constant LR, and no WD.
RMSProp has $\beta_1=0$ and Adam has $\beta_1=.9$; both use
$\beta_2=.95$ and $\epsilon=10^{-8}$.

\paragraph{1. Increase gradient noise.}
I keep the loss, initialization, and optimizers unchanged and replace
the gradient-noise covariance:
\[
\frac{I_2}{B}\quad\longrightarrow\quad\frac{\sigma^2I_2}{B},
\qquad \sigma\in\{1,10,100,300\}.
\]
$I_2$ is the $2\times2$ identity matrix. Increasing $\sigma$ makes
gradient estimates noisier without changing their mean or the number
of updates; $\sigma=1$ is the original setting.
For each noise level and optimizer, I tune LR at $B=1,2,4,8,16,32,64$
and fit $\eta^*(B)=c(B/64)^p$. I freeze the prediction at $B=256$
before separately tuning LR there. The fitted laws are below, with $r=B/64$:

\begin{center}
\small
\begin{tabular}{rll}
\toprule
Noise $\sigma$ & $\eta^*_{\mathrm{RMSProp}}(B)$ & $\eta^*_{\mathrm{Adam}}(B)$ \\
\midrule
1 (original) & $0.025831\,r^{0.90423}$ & $0.025456\,r^{0.90497}$ \\
10  & $0.028842\,r^{0.53622}$ & $0.026957\,r^{0.52044}$ \\
100 & $0.031505\,r^{0.50827}$ & $0.031139\,r^{0.50783}$ \\
300 & $0.019679\,r^{0.50666}$ & $0.019245\,r^{0.50368}$ \\
\bottomrule
\end{tabular}
\end{center}

\paragraph{Takeaway---noise.} Earlier parts found almost linear LR--batch
scaling ($p\approx0.99$ for SGD and $0.90$ for RMSProp/Adam).
As gradient noise increases, the RMSProp/Adam exponent approaches
$0.5$ (Figure~\ref{fig:p31c}, top row). This is square-root scaling,
$\eta^*\propto\sqrt B$: four times the batch size calls for roughly
twice the LR; doubling the batch calls for roughly $1.41$ times the LR.

\paragraph{Intuition for the high-noise result.}
RMSProp rescales each weight's gradient using recent squared gradients.
Ignoring the small numerical-stability constant $\epsilon$, its update
is approximately
\[
\Delta w_i\approx-\frac{\eta g_i}{\sqrt{v_i}},
\]
where $\Delta w_i$ is the change in weight $i$, $g_i$ is its noisy gradient,
and $v_i$ is a running estimate of $\mathbb E[g_i^2]$, the expected
squared gradient, rather than the exact expectation. Adam also smooths
the gradient in the numerator, but uses the same kind of divisor.
Thus, larger $v_i$ reduces the update at a fixed LR.

For our quadratic loss, the true gradient in direction $i$ is $h_iw_i$,
where $h_i$ is the curvature ($h_1=1$, $h_2=10$). Adding independent,
zero-mean gradient noise gives
\[
\mathbb E[g_i^2]=h_i^2\mathbb E[w_i^2]+\sigma^2/B.
\]
The two terms are squared gradient signal and noise variance, respectively.
When noise dominates, $v_i\approx\sigma^2/B$ and
$\sqrt{v_i}\approx\sigma/\sqrt B$. The coefficient multiplying the
gradient is therefore $\eta/\sqrt{v_i}\approx\eta\sqrt B/\sigma$.
At fixed $N$, doubling $B$ halves the number of updates. If this
coefficient grows proportionally to $B$ to compensate, then
$\eta\propto\sqrt B$ at fixed $\sigma$. This is a noise-dominated
heuristic, not an exact law; when signal contributes substantially to
$v_i$, the same argument does not determine the exponent.

\paragraph{2. Replace the loss with a one-parameter problem.}
I repeat the same noise levels, source fits, and batch-256 tests with
$f(w)=w^2/2$, $w_0\sim\mathcal N(0,2)$, and scalar gradient noise
$\epsilon_t\sim\mathcal N(0,\sigma^2/B)$. There is now one weight instead
of two: the direction with curvature ten is removed. Initialization
variance two keeps expected initial loss at one, as before. All other
settings remain unchanged; this comparison changes dimensionality
as well as curvature.

Figure~\ref{fig:p31c} (bottom row) shows the same broad pattern:
$p\approx0.95$--$0.97$ at the original noise level, falling to
$0.49$--$0.51$ at $\sigma\geq100$.

\paragraph{Takeaway---loss change.} The one-parameter loss preserves the
high-noise square-root rule. Noise strength changes the fitted exponent
much more than this loss change, although the low-noise exponents differ
slightly.

\begin{figure}[ht]
    \centering
    \includegraphics[width=\textwidth]{figures/p31c_noise_and_curvature.png}
    \caption{Noise and LR--batch scaling: original two-parameter loss
    (top), one-parameter loss (bottom). Left and middle: RMSProp and Adam
    laws fitted only at $B=1$--$64$ (solid), extrapolated to 256 (dashed).
    Stars are fitted optima; hollow diamonds are frozen predictions at
    256. Right: fitted exponent versus noise; the dotted line is $p=0.5$.
    Largest noise is yellow.}
    \label{fig:p31c}
\end{figure}

\paragraph{Batch-256 check.} At $\sigma\geq100$, predicted LRs give losses
within $1.2\%$ of separately tuned LRs in both settings, although flat
loss curves can hide LR errors. At $\sigma=10$, prediction penalties
are $31\%$/ $24\%$ for RMSProp/Adam with the original loss, versus
$123\%$/ $61\%$ with the one-parameter loss: similarns now. ns now.  scaling trends
do not guarantee similar transfer accuracy.

\paragraph{Takeaway.} More gradient noise moves both RMSProp and Adam toward
$\eta^*\propto\sqrt B$. Replacing the loss with a one-parameter problem
preserves that high-noise rule; noise strength changes the scaling much
more than this loss change, but loss geometry still affects transfer accuracy.

\subsection{(d) Momentum}

At $B=16$ and 256, I tuned SGD LR separately for $\mu=0$ and $.9$.
For Adam I jointly swept LR and 14 values of $\beta_1$, including zero,
$.9$, and values up to $.9999$. I then evaluated the selected pairs on
fresh trajectories; losses below correspond to the curves in
Figure~\ref{fig:p31d}.

\begin{center}
\small
\begin{tabular}{rrrrl}
\toprule
$B$ & SGD, $\mu=0$ & SGD, $\mu=.9$ & Adam, $\beta_1=0$ & Best Adam \\
\midrule
16 & 0.000335 & 0.000307 & 0.000985 & 0.000874 ($\beta_1=.95$) \\
256 & 0.000554 & 0.000914 & 0.012758 & 0.001813 ($\beta_1=.7$) \\
\bottomrule
\end{tabular}
\end{center}

\begin{figure}[ht]
    \centering
    \includegraphics[width=\textwidth]{figures/p31d_momentum.png}
    \caption{Left: final loss after separate LR tuning at each Adam
    momentum setting. Middle/right: learning curves for tuned SGD and
    Adam, with and without momentum. Each curve averages 32,768 fresh
    trajectories; batch 256 permits only 32 updates.}
    \label{fig:p31d}
\end{figure}

\paragraph{SGD.} Fixed momentum $.9$ helps slightly at batch 16 ($9\%$ lower loss),
but hurts at 256 ($65\%$ higher); exact covariance calculations confirm
this finite-budget result. Its small-batch LR is roughly one tenth of
the no-momentum LR, consistent with effective LR $\eta/(1-\mu)$.
One intuition is that momentum combines past gradients, making updates
less sensitive to noisy small-batch gradients. At batch 256, gradients
are already less noisy, so that benefit is smaller. There are also only
32 updates, versus 512 at batch 16: momentum $.9$ remembers roughly ten
updates, so slow adjustment to changing gradients can consume a substantial
fraction of the run. This explains a possible finite-budget trade-off,
not a general rule that momentum hurts large batches.

\paragraph{Adam.} Jointly tuned momentum gives $11\%$ lower loss at batch
16, a modest improvement that is hard to see on the logarithmic learning
curve, versus $86\%$ at 256. The best coefficient decreases from $.95$
to $.7$; values near one perform poorly. Unlike SGD, Adam divides its
gradient estimate by a running estimate of gradient magnitude. Smoothing
the gradient in the numerator can therefore change the behavior of these
normalized updates, not just average away raw gradient noise. A possible
explanation is that moderate momentum reduces oscillatory updates and
permits faster progress: the tuned LR at 256 rises from $.0931$ without
momentum to $.1320$ with $\beta_1=.7$. This is an interpretation, not a
separately isolated mechanism; simply having less gradient noise does
not imply less benefit from Adam momentum.

\paragraph{Takeaway.} Momentum can help more at large batches, but the
coefficient must be tuned: fewer updates can make long gradient memory
harmful. Fixed SGD momentum $.9$ does not reproduce the predicted
large-batch advantage here.

\paragraph{Language-model predictions.}
\begin{itemize}
    \item At small batches, momentum should help smooth noisy gradients,
    though the gain may be modest once LR is well tuned.
    \item At large batches and fixed training tokens, there are fewer
    updates. Well-chosen momentum could help maintain useful directions
    and improve progress, but excessive momentum could slow adaptation.
    \item I would not conclude that SGD momentum generally hurts large
    batches, or that Adam's best coefficients must be $.95$ and $.7$.
    Those results depend on this toy loss and budget.
\end{itemize}

\section{Problem 3.2: Batch size in a language model}
\subsection{(a) Tune LR at fixed weight decay}

I used d8 with WD $.1$, linear LR decay, and 1\% warmup, sweeping seven
peak LRs at each $B\in\{8,16,32,64\}$: nine new runs, 16 reused A1 runs,
and three supplied runs. Reused runs match the model, seeds, schedule,
and seed-42 training prefix. Every run processes exactly
$600{,}000\times1024=614.4$M tokens. Each update uses one microbatch of
$B$ sequences, so doubling $B$ halves the number of updates.

I fit the same quadratic in log LR as in Problem 1 to all seven
measurements at each batch (Figure~\ref{fig:p32a}, left).
The fitted optimum interpolates the sweep;
the best measured loss below comes from an actual run, not the curve.

\begin{center}
\small
\begin{tabular}{rrrrr}
\toprule
$B$ & Updates & Fitted optimal LR & Best measured LR & Best measured loss \\
\midrule
8  & 75,000 & 0.001066 & 0.001  & 2.933064 \\
16 & 37,500 & 0.001803 & 0.0015 & 2.922586 \\
32 & 18,750 & 0.002875 & 0.003  & 2.920011 \\
64 & 9,375  & 0.003547 & 0.003  & 2.925636 \\
\bottomrule
\end{tabular}
\end{center}

\begin{figure}[ht]
    \centering
    \includegraphics[width=\textwidth]{figures/p32a_lr_and_batch.png}
    \caption{Left: measured loss--LR points and quadratic fits in log LR.
    Middle: fitted optima and their power-law fit. Right: best measured
    loss at each batch. Stars denote fitted optima; circles denote actual
    measurements. Batch size controls color, with the largest batch yellow.}
    \label{fig:p32a}
\end{figure}

Fitting a power law to these four optima gives
\[
\eta^*(B)\approx0.003876\left(\frac{B}{64}\right)^{0.588}.
\]
Doubling batch size suggests multiplying peak LR by $1.50$.
Growth also slows between 32 and
64. Batch 32 has the lowest measured loss; at that batch, LR $.003$ and
$.0045$ differ by only $0.000041$ loss. The minima are approximate,
particularly where the curves are flat, and these single-seed sweeps do
not establish statistically significant batch rankings.

\paragraph{Takeaway.} Over batches 8--64, larger batches prefer larger LR,
with close to square-root scaling: $\eta^*\propto B^{0.588}$.
Doubling batch size suggests about $1.50\times$ the LR
($1.41\times$ for exact square-root scaling), not $2\times$.
After LR tuning, loss improves from batch 8 to 32, then
worsens slightly at 64: more updates alone do not guarantee better loss.
Part (b) tests transfer to batches 128 and 256.

\subsection{(b) Compare LR and WD scaling}

I compare two recipes: scale LR with WD fixed at $.1$, or scale WD with
peak LR fixed at $.0015$. The first reuses (a). For the second, I swept
$\lambda\in\{.025,.05,.1,.2,.4,.8,1.6\}$ at each source batch, using
22 new runs and six matching existing runs. All source runs completed
614.4M tokens.

At each batch I fit
$L_B(\lambda)=a_By^2+b_By+c_B$, where $y=\log_2(\lambda/.1)$,
to the seven losses. The optima and measured minima are:

\begin{center}
\small
\begin{tabular}{rrrr}
\toprule
$B$ & Fitted optimal WD & Best measured WD & Best measured loss \\
\midrule
8  & 0.066841 & 0.05 & 2.935132 \\
16 & 0.093868 & 0.1  & 2.922586 \\
32 & 0.133741 & 0.2  & 2.919927 \\
64 & 0.300167 & 0.4  & 2.923157 \\
\bottomrule
\end{tabular}
\end{center}

The power law fitted to these four WD optima
(Figure~\ref{fig:p32b}, middle) is
\[
\lambda^*(B)\approx0.261158\left(\frac{B}{64}\right)^{0.7012}.
\]
It suggests multiplying WD by $1.63$ when doubling batch size; the
best measured WDs instead double exactly. The loss curves are asymmetric:
all-seven-point quadratic fits have loss RMSE $0.008$--$0.009$ and their
minima are sensitive to the high-WD tail. Fitting just three neighboring
points around each measured minimum gives an exponent $0.988$ instead.
Thus the direction is clear, but the precise extrapolation is uncertain.

\begin{figure}[ht]
    \centering
    \includegraphics[width=\textwidth]{figures/p32b_wd_source_fits.png}
    \caption{Left: source WD measurements (circles) and all-seven-point
    quadratic fits in log WD. Middle: fitted optima (stars), their power
    law over source batches (solid), extrapolation (dashed), and frozen
    target predictions (hollow diamonds). Right: measured target losses
    for both recipes; $\Delta L=L_{\mathrm{LR}}-L_{\mathrm{WD}}$.
    Largest batch in each measurement panel is yellow.}
    \label{fig:p32b}
\end{figure}

Both rules were fixed before reading their target results. All four
target runs finished; Figure~\ref{fig:p32b} (right) compares their losses:

\begin{center}
\small
\setlength{\tabcolsep}{4pt}
\begin{tabular}{rrrrrr}
\toprule
$B$ & LR prediction & LR-recipe loss & WD prediction & WD-recipe loss & $\Delta L$ \\
\midrule
128 & 0.005826 & 2.969603 & 0.424593 & 2.947920 & 0.021683 \\
256 & 0.008756 & 3.137107 & 0.690308 & 2.996901 & 0.140206 \\
\bottomrule
\end{tabular}
\end{center}

The LR recipe uses WD $.1$; the WD recipe uses LR $.0015$.
At 128/256, microbatches of 64 are accumulated twice/four times.
Discarding incomplete accumulation groups leaves actual budgets of
614,334,464 and 614,203,392 tokens, respectively, less than $0.04\%$
below nominal.

Neither target has an independent optimum sweep or repeated seeds, so
this comparison ranks the two frozen recipes rather than establishing
optimal hyperparameters or statistical significance.

\paragraph{Takeaways.} Both fitted optimal LR (at fixed WD) and optimal
WD (at fixed LR) increase with batch size. Of the two recipes tested,
fixing LR and scaling WD transfers better than fixing WD and scaling LR:
losses are lower by $.0217$ at batch 128 and $.1402$ at 256.
WD scaling reduces, but does not eliminate, the large-batch loss increase.

\subsection{(c) Ablate momentum}

The experiment fixes the best measured LR--WD pairs: $(.001,.1)$ at
$B=8$ and $(.0015,.690308)$ at $B=256$, with $\beta_2=.95$.
It varies $\beta_1\in\{0,\allowbreak.5,\allowbreak.7,\allowbreak.8,\allowbreak.9,\allowbreak.95,\allowbreak.99\}$, reusing the completed
$.9$ controls. LR and WD are not retuned. These are the best measured
recipes available, not proven joint optima; the batch-256 recipe was
selected from only two transferred recipes. Figure~\ref{fig:p32c}
shows all fourteen completed-budget measurements.
\begin{center}
\begin{tabular}{lrrr}
\toprule
Batch & Loss at $\beta_1=0$ & Best sampled $\beta_1$ & Best loss \\
\midrule
8 & 2.970755 & .95 & 2.929515 \\
256 & 3.172107 & .9 & 2.996901 \\
\bottomrule
\end{tabular}
\end{center}
Relative to no momentum, the best sampled setting lowers loss by $.04124$
at batch 8 and $.17521$ at batch 256. Very long memory is not uniformly
helpful: $\beta_1=.99$ gives loss $3.281432$ at batch 256, worse even than
no momentum. There are approximately 75,000 updates at batch 8 but only
2,343 at batch 256. Momentum averages gradient directions, but increasing
it also delays adaptation. These are one-seed comparisons, not estimates
of statistical uncertainty.
The batch-8 runs process exactly 614,400,000 tokens; the batch-256 runs
use four microbatches of 64 and process 614,203,392 tokens, matching the
completed control's accumulation convention.
\begin{figure}[ht]
\centering
\includegraphics[width=\textwidth]{figures/p32c_momentum.png}
\caption{P3.2(c): measured final loss versus momentum at fixed LR and WD.
Each panel has its own vertical scale so the small-batch differences remain visible.
Circles are measurements; no interpolated momentum optimum is claimed.}
\label{fig:p32c}
\end{figure}
\paragraph{Takeaway.} Momentum helps both batches, especially batch 256,
but too much momentum can hurt. The best sampled coefficients are $.95$
and $.9$, respectively, under these fixed LR--WD recipes.

\subsection{(d) What does the NQM explain?}

\paragraph{LR scaling.} The upward LR trend in Figure~\ref{fig:p32a}
agrees qualitatively with the NQM's competition between cleaner gradients
and fewer updates. Quantitatively, the LM exponent $.588$ differs from
the original-noise Adam NQM's $.905$: doubling batch size suggests
$1.50\times$ versus $1.87\times$ LR. The LM exponent is closer to the
high-noise NQM's $.5$, but this similarity does not establish the same
noise regime in the LM.

\paragraph{Weight decay.} Our 3.1 NQM explicitly has no WD: it predicts
neither the optimal WD exponent nor which LR--WD recipe wins. Thus the
WD recipe's advantage in Figure~\ref{fig:p32b} is not confirmation or
failure of a WD prediction from those simulations.
An additional intuition is that AdamW shrinks weights by a factor
$1-\eta_t\lambda$ per update, where $\eta_t$ is the current LR.
Larger batches give fewer updates and hence fewer decay applications.
At fixed training tokens and schedule shape, preserving approximately
$\eta_{\mathrm{peak}}\lambda/B$ suggests increasing either LR or WD
with batch size. The direction agrees with our measurements, but the
fitted exponents $.588$ and $.701$ are below this simple linear-scaling
hypothesis. This is a shrinkage-preservation argument, not a prediction
measured in our no-WD NQM or a proof of optimal validation loss.

\paragraph{Momentum comparison.} In 3.1(d), tuned Adam momentum
helped much more at the larger batch, while fixed SGD momentum $.9$ hurt
there. LR was retuned for every momentum setting. Part 3.2(c) instead
holds LR and WD fixed. Figure~\ref{fig:p32c} supports the qualitative
prediction that some momentum helps both batches and helps more at the
larger batch. It does not transfer the exact toy optima: the NQM preferred
$.95$ at batch 16 and $.7$ at batch 256; the LM preferred $.95$ at batch 8
and $.9$ at batch 256. The small batches and tuning procedures also differ,
so this is not a controlled quantitative agreement test.

\paragraph{What is missing?} The toy problem has constant curvature,
independent additive gradient noise, and no decoupled WD. The LM has
changing gradient directions, parameter-dependent noise, normalization,
clipping, and a decaying LR schedule. A useful next toy experiment would
add AdamW's actual shrinkage update and repeat the two fixed-LR and
fixed-WD protocols. I would also shift the optimum away from zero, using
$f(w)=\tfrac12(w-w^*)^\top H(w-w^*)$ with $w^*\ne0$:
otherwise decay toward zero is artificially favored by the original toy
loss. With a nonzero optimum, WD can reduce noise while introducing bias,
creating a meaningful tuning trade-off. That would produce a testable
WD prediction; the present no-WD simulations cannot supply one.

\begin{center}
\small
\begin{tabular}{p{.15\textwidth}p{.30\textwidth}p{.43\textwidth}}
\toprule
Prediction & NQM result & LM measurement \\
\midrule
LR versus batch & Adam: LR grows, exponent $.905$ & LR grows, exponent $.588$;
transferred losses $2.969603$ ($B=128$), $3.137107$ ($B=256$). Direction agrees, exponent differs. \\
WD versus batch & Not modeled: no WD in our toy updates & WD-scaling losses
$2.947920$, $2.996901$, better than LR scaling by $.021683$, $.140206$. No toy WD prediction to falsify. \\
Momentum & Tuned Adam helps both, more at larger batch & Fixed-LR LM also improves both:
loss drops $.041240$ at $B=8$, $.175206$ at $B=256$. Best coefficients do not transfer exactly. \\
\bottomrule
\end{tabular}
\end{center}

\paragraph{Takeaway.} The NQM gets the direction of LR scaling and the
larger-batch momentum benefit right, but not the exact exponents or best
coefficients. Our no-WD version cannot explain which WD recipe wins.

\section{Problem 4.0: Alignment-dependent scaling}

The goal is to keep features usable and updates nonvanishing as model
width grows, so the same base LR can potentially transfer across widths.
For a linear layer, the handout writes
\[
\begin{aligned}
z&=n^{-a}Wx, & \operatorname{Var}(W_{0,ij})&\propto n^{-2b},\\
\eta_W&=\eta n^{-c}, & \Delta W&=-\eta n^{-c}U.
\end{aligned}
\]
Here $a$ controls the forward multiplier, $b$ initialization standard
deviation (hence variance uses $2b$), and $c$ LR. The base LR $\eta$ is
width independent; the normalized Adam update direction $U$ is assumed
to have order-one RMS. Initial weights have independent centered
Gaussian entries and are independent of the initial inputs.
RMS means the square root of the mean squared
entries; order one means bounded and nonvanishing with width.
All scales are evaluated at a fixed training step as width grows.
Writing a scale as $\Theta(n^r)$ means its typical size is proportional
to $n^r$, up to width-independent factors. A positive $r$ means it grows
with width, a negative $r$ means it vanishes, and $r=0$ means it stays
order one. Our derivation will therefore set the relevant exponents to zero.

There are three layer types. Input matrices map a fixed-size external
input into width-$n$ features, so their fan-in $d$ stays fixed.
Hidden matrices map between features whose sizes both grow with $n$.
The readout maps width-$n$ features to a fixed number $q$ of output logits.
Let $x_0$ be the pre-update input and $\Delta x=x_1-x_0$ its change,
both with order-one RMS for hidden and readout matrices.

\paragraph{The two alignment choices.}
First, how strongly does the update direction $U$ act on its current
input $x_0$? Second, how strongly do the initial readout weights act on
the change $\Delta x$ in their input? These are different relationships.
The ratios below divide the size of the resulting output by the RMS
sizes of the matrix and input, leaving their amplification due to alignment:
\[
\begin{aligned}
R(U,x_0)&=\frac{\|Ux_0\|_{\rm RMS}}
{\|U\|_{\rm RMS}\|x_0\|_{\rm RMS}}=\Theta(n^\alpha),\\
S(V_0,\Delta x)&=\frac{\|V_0^\top\Delta x\|_{\rm RMS}}
{\|V_0\|_{\rm RMS}\|\Delta x\|_{\rm RMS}}=\Theta(n^\omega).
\end{aligned}
\]
where $V_0=W_0^\top$ for the readout and
$\omega=\omega_{\rm move}$. For each ratio independently, full alignment
means exponent $1$ and no alignment means $1/2$. The intuition is that
aligned contributions add coherently, giving amplification proportional
to $n$; independent centered contributions partly cancel, giving
amplification proportional to $\sqrt n$. We choose each assumption
separately, producing four combinations.

\subsection{(a) Derive the four combinations}

\paragraph{Step 1: keep initial hidden features order one.}
Each output coordinate sums $n$ initially independent weight contributions.
Because cross terms have zero expectation,
\[
\mathbb E[z_{0,i}^2\mid x_0]
\propto n^{-2(a+b)}\sum_{j=1}^n x_{0,j}^2
=n^{1-2(a+b)}\|x_0\|_{\rm RMS}^2.
\]
Taking a square root gives initial RMS scale $n^{1/2-a-b}$.
Setting its exponent to zero requires $a+b=1/2$. The question fixes
$a=0$ for hidden matrices, so $b=1/2$. This does not depend on either
training-alignment assumption: initialization comes before any update.

\paragraph{Step 2: keep the direct weight update order one.}
The forward multiplier contributes $n^{-a}$, the update magnitude
contributes $\eta n^{-c}$, and update alignment contributes $n^\alpha$:
\[
\|n^{-a}\Delta W x_0\|_{\rm RMS}
=\Theta(\eta n^{\alpha-a-c}).
\]
Thus $a+c=\alpha$. This applies to hidden and readout matrices.
For hidden matrices, $a=0$ gives $c=\alpha$; together with Step 1,
their tuple is $(0,1/2,\alpha)$.
For the readout, the question instead fixes $c=0$, so $a=\alpha$.

\paragraph{Step 3: keep the readout's response to feature movement order one.}
Even before its own weights change, the initial readout sees a changed
input $\Delta x$. Its forward multiplier is $n^{-a}$, its stored weights
have RMS scale $n^{-b}$, and this second alignment contributes $n^\omega$:
\[
\|n^{-a}V_0^\top\Delta x\|_{\rm RMS}
=\Theta(n^{\omega-a-b}).
\]
Setting this exponent to zero requires $a+b=\omega$. Step 2 already
gave $a=\alpha$, so $b=\omega-\alpha$. The readout tuple is therefore
$(\alpha,\omega-\alpha,0)$.
Unlike hidden features, initial readout logits may vanish; they need
only remain bounded. Their scale is
$\Theta(n^{1/2-\omega})$, constant for no readout alignment and vanishing
as $n^{-1/2}$ for full readout alignment. Initial weights are independent
of $x_0$, even though they may align with the learned $\Delta x$.

\paragraph{Step 4: fixed-fan-in input matrices.}
Here we sum over a fixed $d$, not a growing $n$, so neither initialization
nor update amplification introduces a width factor. Their scales are
$n^{-a-b}$ and $\eta n^{-a-c}$. Setting both exponents to zero gives
$a+b=a+c=0$. With the required $a=0$, the input tuple is $(0,0,0)$
in every case; the external input also has no $\Delta x$.

\paragraph{Step 5: substitute each pair of assumptions.}
We now plug the four pairs into hidden $(0,1/2,\alpha)$ and readout
$(\alpha,\omega-\alpha,0)$:
\begin{enumerate}
    \item \textbf{Full update, full readout alignment:}
    $\alpha=1$, $\omega=1$. Hidden $c=1$; readout $a=1$ and
    $b=1-1=0$. The tuples are $(0,1/2,1)$ and $(1,0,0)$.
    \item \textbf{Full update, no readout alignment:}
    $\alpha=1$, $\omega=1/2$. Hidden $c=1$; readout $a=1$ and
    $b=1/2-1=-1/2$. The tuples are $(0,1/2,1)$ and $(1,-1/2,0)$.
    \item \textbf{No update, full readout alignment:}
    $\alpha=1/2$, $\omega=1$. Hidden $c=1/2$; readout $a=1/2$ and
    $b=1-1/2=1/2$. The tuples are $(0,1/2,1/2)$ and $(1/2,1/2,0)$.
    \item \textbf{No update, no readout alignment:}
    $\alpha=1/2$, $\omega=1/2$. Hidden $c=1/2$; readout $a=1/2$ and
    $b=1/2-1/2=0$. The tuples are $(0,1/2,1/2)$ and $(1/2,0,0)$.
\end{enumerate}
Input tuples remain $(0,0,0)$. Table~\ref{tab:p40-abc} summarizes these
substitutions.

\begin{table}[ht]
\centering
\small
\begin{tabular}{llccc}
\toprule
\shortstack{Update\\alignment} & \shortstack{Readout\\alignment} &
\shortstack{Hidden\\$(a,b,c)$} & \shortstack{Readout\\$(a,b,c)$} & \shortstack{Input\\$(a,b,c)$} \\
\midrule
Full & Full & $(0,\frac12,1)$ & $(1,0,0)$ & $(0,0,0)$ \\
Full & None & $(0,\frac12,1)$ & $(1,-\frac12,0)$ & $(0,0,0)$ \\
None & Full & $(0,\frac12,\frac12)$ & $(\frac12,\frac12,0)$ & $(0,0,0)$ \\
None & None & $(0,\frac12,\frac12)$ & $(\frac12,0,0)$ & $(0,0,0)$ \\
\bottomrule
\end{tabular}
\caption{Exponent choices under the handout's two independent alignment assumptions.}
\label{tab:p40-abc}
\end{table}

\paragraph{Takeaway.} The two assumptions have different jobs: update
alignment sets how strongly to scale down direct updates ($a+c$), while
readout--feature alignment sets the head's initialization and forward
scaling ($a+b$). The four cases come from combining these two choices.

\subsection{(b) Forward multiplier, initialization variance, and LR}

Let $m=n/n_0$, with reference width $n_0$, and let $k$ be each hidden
matrix's actual fan-in. A forward multiplier scales as $m^{-a}$,
initialization variance as $m^{-2b}$, and LR as $\eta m^{-c}$.
For example, the full/full readout has $a=1$, $b=0$, $c=0$:
its forward multiplier is $1/m$, its initialization variance does not change,
and its LR does not change. The full/none case instead has $b=-1/2$,
so its variance multiplier is $m^{-2(-1/2)}=m$.
To match the handout's preceding table at $n_0$, the reference
initialization variances are $1/d$ for input, $1/k_0$ for hidden, and
$1/n_0$ for readout, where $k_0$ is the hidden matrix's fan-in at $n_0$.
Because hidden fan-in scales with width, $k=mk_0$, its variance becomes
$(1/k_0)m^{-1}=1/k$.
Table~\ref{tab:p40-rules} gives all three requested quantities for
every layer type and alignment combination. Case names list update
alignment first and initial-readout/feature-change alignment second.

\begin{table}[ht]
\centering
\small
\begin{tabular}{llccc}
\toprule
Alignment case & Layer & \shortstack{Forward\\multiplier} &
\shortstack{Initialization\\variance} & LR \\
\midrule
Full / full & Input & $1$ & $1/d$ & $\eta$ \\
 & Hidden & $1$ & $1/k$ & $\eta/m$ \\
 & Readout & $1/m$ & $1/n_0$ & $\eta$ \\
\midrule
Full / none & Input & $1$ & $1/d$ & $\eta$ \\
 & Hidden & $1$ & $1/k$ & $\eta/m$ \\
 & Readout & $1/m$ & $m/n_0$ & $\eta$ \\
\midrule
None / full & Input & $1$ & $1/d$ & $\eta$ \\
 & Hidden & $1$ & $1/k$ & $\eta/\sqrt m$ \\
 & Readout & $1/\sqrt m$ & $1/(mn_0)$ & $\eta$ \\
\midrule
None / none & Input & $1$ & $1/d$ & $\eta$ \\
 & Hidden & $1$ & $1/k$ & $\eta/\sqrt m$ \\
 & Readout & $1/\sqrt m$ & $1/n_0$ & $\eta$ \\
\bottomrule
\end{tabular}
\caption{Forward multiplier, initialization variance, and LR for each
layer type. At $n=n_0$ ($m=1$), every case matches the handout's preceding table.}
\label{tab:p40-rules}
\end{table}

At $m=1$, all forward multipliers are $1$, all LRs are $\eta$, and
initialization variances are exactly $1/d$, $1/k_0$, and $1/n_0$.

Removing update alignment changes hidden LR from $\eta/m$ to
$\eta/\sqrt m$ and the readout multiplier from $1/m$ to $1/\sqrt m$.
Removing readout alignment instead multiplies its initialization variance
by $m$ at fixed update alignment. The negative $b$ in the full/none case
is intentional: readout initialization variance grows, but its forward multiplier
keeps initial logits bounded.

\paragraph{Interaction check.} The total change is
\[
\Delta z=n^{-a}(\Delta W x_0+W_0\Delta x+\Delta W\Delta x).
\]
The handout additionally assumes $R(U,\Delta x)=O(n^\alpha)$, giving
\[
\|n^{-a}\Delta W\Delta x\|_{\rm RMS}
=O(\eta n^{\alpha-a-c})=O(\eta)
\]
for both hidden and readout matrices in all four cases.
Without this extra assumption, a general bound proportional to
$n^{1-a-c}$ would not guarantee bounded interaction in the no-update-
alignment cases. For hidden matrices, the Gaussian operator-norm bound
(controlling the largest possible vector amplification by $W_0$)
also gives $\|n^{-a}W_0\Delta x\|_{\rm RMS}=O(n^{1/2-a-b})=O(1)$,
even for feature changes dependent on $W_0$. For the readout, this term
is order one by $a+b=\omega$. Input matrices have $\Delta x=0$.
Thus all three terms stay bounded under the stated assumptions.

\paragraph{Takeaway.} Full alignment in both places recovers the
handout's $\mu$P rules: hidden LR $\eta/m$, readout multiplier $1/m$,
and readout variance $1/n_0$. Other alignment assumptions give different
width rules; these are conditional predictions, not evidence that a real
Transformer satisfies the assumptions or transfers LR successfully.

One finite-width caution matters for the experiments below. With readout
$W\in\mathbb R^{q\times n}$, pooling RMS over all $q$ vocabulary outputs
gives the exact bound
\[
 S(W,h)\le\frac{n\,\|W\|_{\rm op}}{\|W\|_F},
\]
where $\|W\|_{\rm op}$ is its largest possible vector amplification and
$\|W\|_F$ is the square root of its summed squared entries. A Gaussian
readout with many more outputs than inputs is approximately an isometry
(see \href{https://arxiv.org/abs/1011.3027}{Vershynin, Corollary 5.35}):
it gives ratios near $\sqrt n$ for many directions, including directions
that depend on it. Here $q=4{,}096$, especially large relative to the
depth experiment's width 64. Thus an observed $\omega_{\rm move}\approx.5$
does not by itself prove independence or absence of feature learning;
the pooled ratio is constrained by the initial matrix's finite-width
geometry. The $.5$ and $1$ lines remain useful reference scales, not
binary labels for every measured update.

\clearpage
\section{Problem 4.1: Five-update Transformer stress test}

\paragraph{What is being transferred?} The \emph{base learning rate}
$\eta$ is one tunable number. A parameterization converts it into actual
per-layer rates. For example, at width $2{,}560$, $m=n/512=5$, so
$\mu$P's hidden matrices use $\eta/5$, while embeddings, readout, and
normalization gains still use $\eta$. Successful transfer means that
approximately the same \emph{base} LR works at different dimensions;
it does not mean every matrix receives the same actual LR.

\paragraph{Controlled protocol.} The policies are implemented in
\path{experiments/a2/p41_policies.py}, using the supplied stress architecture
and measurement code. Every configuration uses the same five batches of 64
length-1,024 sequences, the same 64 validation sequences, seed 42,
constant LR, no warmup, no WD, and no clipping. File hashes verify the
shared data: the first 320 contiguous training sequences and first 64
validation sequences of the shared split. Diagnostic probes use one fixed validation sequence, pooling
all 1,024 token positions; validation loss uses all 64 sequences.
Microbatches only accumulate gradients to the same full batch.
All parameters are initialized explicitly; none retain the scaffold's NaNs.

\subsection{(a) Width: transfer and tuned loss}

Depth is two, head dimension is 64, and reference width is 512.
The width experiment uses FP32 weights, activations, gradients and Adam
states, math attention and disabled TF32. Embedding variance is one,
hidden-matrix variance is $1/k$ for actual fan-in $k$, and norm gains start
at one. Kaiming uses readout variance $1/n$, multiplier one and hidden LR
$\eta$. Standard $\mu$P instead uses variance $1/512$, multiplier $512/n$
and hidden LR $512\eta/n$. Other parameter rates remain $\eta$.

Before the full sweeps I predicted that $\mu$P would transfer its base LR
more reliably, but did not predict that it must achieve lower tuned loss.
I sweep widths $640$, $2{,}560$ and $5{,}120$. The initial Kaiming grid is
$\{.0003,.001,.003,.01,.03\}$; the $\mu$P grid is
$\{.001,.003,.01,.03,.1\}$. Timing pilots are reused, not repeated.
An endpoint minimum requires an additional bracketing LR before an
interpolated optimum is reported.

\paragraph{Curve fitting.} To estimate each optimum, I interpolate the
best measured point and its nearest lower/higher LRs with
\[
 L(\eta)=a x^2+b x+c,\qquad x=\log_2(\eta/.003),\qquad
 \eta^*=.003\,2^{-b/(2a)}.
\]
I require $a>0$ and an optimum inside those neighbors. This local fit
avoids letting a high-loss tail dominate the minimum; a quadratic over
the entire grid is retained as a sensitivity check. Three points determine
a quadratic exactly, not an uncertainty estimate. Fitted minimum loss
and best measured loss are reported separately.

The completed width sweep contains 38 unique configurations: 30 initial
points and eight lower-LR bracketing points. Figure~\ref{fig:p41-width-fits}
shows the results:
\begin{center}
\begin{tabular}{llrrr}
\toprule
Policy & Width & Fitted LR & Fitted loss & Best measured loss \\
\midrule
Kaiming & 640 & .001023 & 7.0686 & 7.0688 \\
Kaiming & 2,560 & .000146 & 7.2700 & 7.3131 \\
Kaiming & 5,120 & .0000454 & 7.4882 & 7.5375 \\
$\mu$P & 640 & .001269 & 6.8822 & 6.9094 \\
$\mu$P & 2,560 & .001155 & 6.7367 & 6.7464 \\
$\mu$P & 5,120 & .001079 & 6.7417 & 6.7445 \\
\bottomrule
\end{tabular}
\end{center}
Across an $8\times$ width increase, Kaiming's fitted optimal LR falls
$22.6\times$; $\mu$P's optima span only $1.18\times$. The best sampled
$\mu$P LR is $.001$ at every width. Kaiming instead selects $.001$,
$.0001$ and $.00003$. Transferring its width-640 rate $.001$ costs
$1.2575$ loss at width 2,560 and $15.2529$ at 5,120, relative to their
best sampled rates. $\mu$P incurs no sampled-grid transfer penalty.
It also has lower best measured loss at every width, by $.1594$, $.5667$
and $.7931$, respectively. These conclusions concern this five-update,
single-seed protocol, not arbitrary training lengths.

\begin{figure}[htbp]
\centering
\includegraphics[width=\linewidth]{figures/p41_width_fits.png}
\caption{Five-update width test: measured losses (circles), local
log-LR fits and fitted optima (stars). Main LR panels zoom into the
basins; insets retain the full measurement range, including high-loss
tails. Bottom: optimal base LR and fitted minimum loss versus width.
The largest width is yellow.}
\label{fig:p41-width-fits}
\end{figure}

\paragraph{Takeaway.} In the five-update test, $\mu$P transfers the same
base LR across an $8\times$ width range and also achieves lower tuned
loss. Kaiming needs a much smaller LR as width grows.

\subsection{(b) Width diagnostics}

For each best \emph{sampled} LR, I track logit RMS, final-RMSNorm output
RMS and movement from initialization through all five updates. The norm
output is measured before the readout multiplier. For every Q/K/V/O,
FF gate/up/down matrix and readout, the diagnostic uses the actual matrix
change and its pre-update input:
\[
 \alpha_{\rm upd}^{(t)}=\log_k\frac{\|\Delta A_tu_{t-1}\|_{\rm RMS}}
 {\|\Delta A_t\|_{\rm RMS}\|u_{t-1}\|_{\rm RMS}}.
\]
Values near $.5$ resemble the random, unaligned reference; values near
$1$ resemble full alignment. The second diagnostic
$\omega_{\rm move}=\log_n S(V_0,h_t-h_0)$ uses the \emph{initial}
readout and change in normalized features, rather than the trained
readout or current features. It is omitted at initialization when movement
is zero. These finite-width measurements test assumptions, not exact
asymptotic exponents.

Figure~\ref{fig:p41-width-features} separates feature magnitude from
feature change. Both policies keep the final normalized RMS close to one,
but their movements differ. At update five, Kaiming movement falls from
$1.029$ to $.524$ to $.321$ as width increases; $\mu$P retains
$1.038$, $1.007$ and $.954$. Thus $\mu$P preserves substantial feature
learning across these widths, while Kaiming's retuned larger models
change their features less. Three finite widths do not establish a
vanishing infinite-width limit. Moreover, aggregate movement includes
embedding, normalization and hidden-weight effects; it does not isolate
every layer's contribution.

$\mu$P's initial logit RMS decreases from $.893$ to $.447$ to $.316$,
as expected from its readout multiplier. By update five the values are
$1.156$, $.986$ and $.889$: the initially smaller logits do not prevent
learning. Kaiming's retuned logit RMS stays between $1.016$ and $1.237$.

\begin{figure}[htbp]
\centering
\includegraphics[width=\linewidth]{figures/p41_width_features.png}
\caption{Best sampled width configurations through five updates.
Normalized RMS stays near one under both policies (the second column
zooms into this small range), but only $\mu$P keeps movement approximately
width-independent. The last column uses the initial readout, with $.5$
and $1$ reference lines; zero initial movement has no defined exponent.}
\label{fig:p41-width-features}
\end{figure}

Figure~\ref{fig:p41-width-alignment} shows that a single alignment
assumption does not describe every matrix. At update five, $\mu$P's
median update alignment is $.678$, $.734$ and $.757$ across the three
widths. Attention-output projections and the readout are more strongly
aligned (projection medians about $.907$--$.930$), particularly in the
later block; query/key and FF projections are generally less aligned.
Kaiming's overall medians are $.725$, $.714$ and $.720$.
These measurements support substantial, nonuniform update alignment,
not exact full alignment everywhere. Readout--movement exponents instead
stay near $.5$ ($.501$--$.529$ at update five). As discussed in 4.0,
the finite readout geometry constrains this diagnostic, so it cannot alone
establish independent features or disprove feature learning.

\begin{figure}[htbp]
\centering
\includegraphics[width=\linewidth]{figures/p41_width_alignment.png}
\caption{Actual-update alignment for both blocks' Q/K/V/O and FF
gate/up/down matrices, plus readout, at the best sampled LRs.
Every matrix uses its pre-update probe input and actual fan-in.
Attention-output and readout updates are more aligned than most other
matrices; alignment also differs between blocks and training steps.}
\label{fig:p41-width-alignment}
\end{figure}

\paragraph{Takeaway.} Unit normalized RMS is not enough: $\mu$P also
preserves feature movement. Update alignment varies by matrix, and a
readout ratio near $.5$ does not mean that features stopped learning.

\subsection{(c) Depth: transfer and tuned loss}

Width and reference width are both 64, head dimension is 64, and reference
depth is two. Parameters, gradients, optimizer states and residual sums
remain FP32; branches use BF16 autocast. Write $r=L/2$. The policies are:
\begin{center}
\begin{tabular}{lccc}
\toprule
Policy & Branch multiplier & Block LR & Block Adam $\epsilon$ \\
\midrule
Standard $\mu$P & $1$ & $\eta$ & $10^{-8}$ \\
\href{https://arxiv.org/abs/2310.02244}{Depth-$\mu$P} & $r^{-1/2}$ & $\eta r^{-1/2}$ & $10^{-8}r^{-1/2}$ \\
\href{https://arxiv.org/abs/2505.01618}{CompleteP} & $r^{-1}$ & $\eta$ & $10^{-8}r^{-1}$ \\
\bottomrule
\end{tabular}
\end{center}
``Block'' includes hidden matrices and all within-block norms, including
Q/K normalization. Embedding, readout and final norm keep their original
LR and Adam stabilizer. I test depths $2$, $100$, $1{,}000$, initially at
$\{.001,.003,.01,.03,.1\}$. At depth two the policies coincide, so one
measured curve is reused for all three. I predicted that stronger
residual damping would improve transfer, especially CompleteP, but kept
the tuned-loss comparison separate from that prediction.
A common $.02$ measurement resolves the wide $.01$--$.03$ bracket before
comparing interpolated minima; it is added to all seven unique configurations.

\paragraph{Why these multipliers?} As a heuristic, independent branch
contributions add like $\sqrt r$, while coherent changes can add like $r$.
Depth-$\mu$P damps each branch by $r^{-1/2}$ and each block's LR by
$r^{-1/2}$: their product is $r^{-1}$. CompleteP instead puts the whole
$r^{-1}$ factor on the branch, leaving block LR unchanged. Thus both
can control accumulated changes without necessarily making feature
movement vanish. Transformer branches are not independent, so this
intuition is not a proof; the next probes test it.
Adam's stabilizer is scaled too because damping a gradient by $s$ scales
its first moment and square-root second moment by $s$. Replacing
$\epsilon$ with $s\epsilon$ keeps its relative importance unchanged.

All 42 unique depth configurations are complete, including the common
$.02$ refinement. Figure~\ref{fig:p41-depth-fits} shows the local fits;
depth two supplies the same reference curve to all policies.
\begin{center}
\begin{tabular}{llrrr}
\toprule
Policy & Depth & Fitted LR & Fitted loss & Best measured loss \\
\midrule
All & 2 & .01738 & 6.8319 & 6.8462 \\
Standard $\mu$P & 100 & .01865 & 7.0764 & 7.0774 \\
Standard $\mu$P & 1,000 & .01747 & 7.0694 & 7.0771 \\
Depth-$\mu$P & 100 & .01819 & 6.8155 & 6.8225 \\
Depth-$\mu$P & 1,000 & .01756 & 6.9757 & 6.9786 \\
CompleteP & 100 & .01682 & 6.8039 & 6.8279 \\
CompleteP & 1,000 & .01605 & 6.8538 & 6.8999 \\
\bottomrule
\end{tabular}
\end{center}
Every configuration selects the same sampled LR $.02$. The fitted LR
ranges are also narrow: about $7\%$ for standard $\mu$P, $5\%$ for
Depth-$\mu$P and $8\%$ for CompleteP. This does not establish CompleteP
as a better LR-transfer scheme in this test, contrary to my initial
prediction. Yet transfer and loss clearly separate: at depth 1,000,
CompleteP's measured loss is $.1772$ below standard $\mu$P and $.0787$
below Depth-$\mu$P. At depth 100, Depth-$\mu$P wins the measured comparison
over CompleteP by only $.0055$, while the fitted minima reverse this
small ordering. I would not infer a reliable winner there without repeats.

\begin{figure}[htbp]
\centering
\includegraphics[width=\linewidth]{figures/p41_depth_fits.png}
\caption{Five-update depth test: loss measurements, local log-LR fits
and fitted optima; depth two is reused across policies. Bottom panels
separate LR transfer from fitted minimum loss. All best sampled rates
are $.02$, and the fitted-rate differences are small. Largest depth is yellow.}
\label{fig:p41-depth-fits}
\end{figure}

\paragraph{Takeaway.} All three policies transfer approximately the same
base LR here. Depth corrections improve deep-model loss, but lower loss
does not imply a noticeably more invariant optimal LR.

\subsection{(d) Depth diagnostics}

At best sampled rates, I compare pre-normalization residual RMS, branch
RMS \emph{before} its depth multiplier, movement of normalized features
at the midpoint and final positions, and $\omega_{\rm move}$ over five updates.
RMSNorm can keep its output near unit scale even when its input residual
stream grows. Consequently, normalized output RMS alone is not a
sufficient stability check.

The completed depth-1,000 sweeps give a concrete example. All three
policies select sampled LR $.02$; after five updates:
\begin{center}
\begin{tabular}{lrrrr}
\toprule
Policy & Residual RMS & Final norm RMS & Movement & $\omega_{\rm move}$ \\
\midrule
Standard $\mu$P & 10,219.58 & .961 & 1.047 & .5000 \\
Depth-$\mu$P & 8.72 & .952 & 1.238 & .4980 \\
CompleteP & 7.95 & .942 & 1.264 & .4994 \\
\bottomrule
\end{tabular}
\end{center}
The normalized scales look similar despite over $1{,}000\times$ larger
residual magnitude without a depth correction. Both corrections retain
order-one normalized-feature movement rather than simply freezing the
features. Nearly unchanged normalized RMS does not establish controlled
residual growth. We also need to check the raw stream.

Figure~\ref{fig:p41-depth-signals} shows how this happens. Between depths
100 and 1,000, standard $\mu$P's final residual grows $9.92\times$.
Depth-$\mu$P's grows only $1.75\times$, while CompleteP's remains almost
unchanged ($7.98$ versus $7.95$). The individual unscaled branches need
not explode for their sum to grow: standard $\mu$P's last FF branch RMS
is $12.89$ versus $13.32$ at those depths, despite the much larger
residual. The branch multiplier controls accumulation across layers,
not just the magnitude of a single branch.

Figure~\ref{fig:p41-depth-features} checks movement at the midpoint and
final normalized positions. CompleteP's final movement stays near
$1.24$--$1.26$ over all depths, and Depth-$\mu$P's near $1.18$--$1.24$.
At the midpoint, movement also stays substantial; neither policy merely
suppresses all feature changes. All readout--movement ratios remain close
to $.5$, subject to the finite-geometry caveat in 4.0.
Thus the corrected policies strongly reduce residual accumulation while
retaining feature learning on these probes. Five updates and three depths
cannot prove asymptotic boundedness. Conversely, standard $\mu$P's good
LR transfer in (c) does not certify controlled internal signals.

\begin{figure}[htbp]
\centering
\includegraphics[width=\linewidth]{figures/p41_depth_signals.png}
\caption{Best sampled depth runs: residual RMS before final normalization,
and FF/attention branch RMS \emph{before} depth multiplication. Branch
panels use solid lines for the last block and dashed lines for the midpoint.
Standard $\mu$P accumulates a large stream even without enormous individual
branches. Vertical scales differ across policies; all RMS panels are logarithmic.}
\label{fig:p41-depth-signals}
\end{figure}
\begin{figure}[htbp]
\centering
\includegraphics[width=.88\linewidth]{figures/p41_depth_features.png}
\caption{Normalized-feature movement at final (solid) and midpoint
(dashed) positions, and initial-readout alignment with final feature
change. Depth corrections preserve nonzero movement while controlling
the raw stream in Figure~\ref{fig:p41-depth-signals}. Largest depth is yellow.}
\label{fig:p41-depth-features}
\end{figure}

\paragraph{Takeaway.} A good transferred LR and unit normalized outputs
can hide a growing residual stream. Depth corrections control that
accumulation without freezing feature changes; CompleteP gives the most
depth-independent residual magnitude in this test.

\clearpage
\section{Problem 4.2: Longer-training transfer}

\paragraph{Why repeat the experiment?} Five updates test early feature
and update scales. They do not establish the best recipe after learning
from many tokens, with warmup, LR decay, WD and clipping. This experiment
uses the default LM at 153.6M tokens, with those settings restored.
Its initialization is intentionally different from the stress test:
the course baseline has small, shared embedding/readout draws, not
embedding variance one.

The longer runs have 2,344 optimizer steps and 23 warmup steps. The first
step uses LR zero, so feature movement is zero and update alignment is
undefined; the first five checkpoints are still within warmup.
They are therefore not the same intervention as 4.1's five full-rate
updates from initialization.

\paragraph{Implementation and controls.} For a shared standard-Gaussian
matrix $G$ truncated at $\pm3$, the course baseline uses embedding $G/n$,
readout $G/\sqrt n$, readout multiplier one and hidden LR $\eta$.
$\mu$P uses embedding $G/512$, readout $G/\sqrt{512}$, multiplier $512/n$
and hidden LR $512\eta/n$. These are the only four width-policy changes.
Sharing the initial draws does not tie the weights: embedding and readout
remain separate trainable matrices.
Both prescriptions use the same initialization draws, ordered seed-42
training prefix, architecture, norm $\epsilon=10^{-5}$, clipping at one,
WD $.1$, Adam betas $(.9,.95)$ and linear schedule with 1\% warmup.
In a CPU check at width 512 I verified exact agreement with the default model's weights,
logits and an AdamW update. The supplied width-512 runs and diagnostics
are reused rather than retrained.
Feature and alignment diagnostics pool the first eight fixed validation
sequences (including every token position); final validation loss uses
all 1,000 sequences. Thus the probe inputs stay fixed within and across
these comparisons, while the loss is measured on the full validation set.

\subsection{(a) Transfer the reference LR}

Problem 1(a)'s best sampled LR at 153.6M tokens is $.003$.
I test both prescriptions at widths 128 and 256 with an equal initial grid
$\{.00075,\allowbreak.0015,\allowbreak.003,\allowbreak.006,\allowbreak.012\}$, including direct $.003$ transfer.
Before examining the full source sweeps, I predicted better base-LR
transfer with $\mu$P, with possible shifts from clipping, norm epsilon
and WD. Update alignment is measured for all matrices throughout
training, not just during the first five updates.

At width 128, direct $.003$ transfer gives loss $3.6940$ in the baseline
and $3.6162$ in $\mu$P. Their best sampled losses are $3.6274$ and
$3.6162$: transfer costs $.0666$ for the baseline but nothing on this
grid for $\mu$P. At width 256, the corresponding gaps are only $.0006$
and $.0063$. Thus $\mu$P transfers much better at width 128, but does
not win every individual comparison (Figure~\ref{fig:p42-source-fits}).

Figures~\ref{fig:p42-align-baseline} and \ref{fig:p42-align-mup} show
each matrix's update alignment at the transferred rate. The median
over 57 matrices changes as follows; these are descriptive medians,
not fitted width exponents or confidence intervals.
\begin{center}
\begin{tabular}{llrrr}
\toprule
Prescription & Width & Update 2 & Update 5 & Final update \\
\midrule
Baseline & 128 & .673 & .799 & .613 \\
Baseline & 256 & .746 & .905 & .611 \\
$\mu$P & 128 & .668 & .879 & .597 \\
$\mu$P & 256 & .731 & .910 & .611 \\
Both (reference) & 512 & .746 & .919 & .624 \\
\bottomrule
\end{tabular}
\end{center}

\begin{figure}[htbp]
\centering
\includegraphics[width=\linewidth]{figures/p42_source_alignment_baseline.png}
\caption{Baseline update alignment at direct LR $.003$ for every attention,
FF and readout matrix. Rows are equally spaced \emph{recorded checkpoints},
not equal time intervals; labels give actual update numbers. The dashed
boundary separates the first five updates from later training. White at
update one denotes the zero warmup update, for which alignment is undefined.}
\label{fig:p42-align-baseline}
\end{figure}
\begin{figure}[htbp]
\centering
\includegraphics[width=\linewidth]{figures/p42_source_alignment_mup.png}
\caption{$\mu$P update alignment, using the same matrices, source rate,
color scale and checkpoint convention as Figure~\ref{fig:p42-align-baseline}.
Strong early alignment weakens later, with substantial differences between
matrix types. FF down projections use their actual fan-in $3.5n$.}
\label{fig:p42-align-mup}
\end{figure}

In the supplied width-512 reference, median update alignment over all
57 matrices rises from $.746$ at update two to $.919$ at update five,
then falls to $.624$ at update 2,344. Alignment is therefore stage-dependent,
not one fixed full/no-alignment choice throughout training.
Readout--movement alignment instead stays near $.5$: $.4976$ at update
five and $.5007$ at the end, while normalized-feature movement remains
nonzero ($1.288$ and $1.335$). A near-random readout ratio does not imply
vanishing feature changes.
Figure~\ref{fig:p42-source-omega} shows the same near-$.5$ behavior at
all three widths under both prescriptions.

\begin{figure}[htbp]
\centering
\includegraphics[width=.9\linewidth]{figures/p42_source_readout_alignment.png}
\caption{Readout--normalized-feature-change alignment at the transferred
LR. The horizontal lines are the $.5$ and $1$ reference scales; zero
movement at initialization/update one has no defined exponent.}
\label{fig:p42-source-omega}
\end{figure}

\paragraph{Takeaway.} $\mu$P makes the reference LR a much better recipe
for width 128. Update alignment is strong early and weaker late, while
readout--movement alignment stays near $.5$ throughout.

\subsection{(b) Transfer versus tuned performance}

The source curves use the same bracketed log-LR quadratic interpolation
as 4.1. I compare fitted optimal base LR, fitted minimum validation loss,
and the measured loss at $.003$. Better LR transfer and lower achievable
loss are distinct questions.

\begin{center}
\begin{tabular}{llrrrr}
\toprule
Policy & Width & Fitted LR & Fitted loss & Best measured & Loss at $.003$ \\
\midrule
Baseline & 128 & .008283 & 3.6200 & 3.6274 & 3.6940 \\
Baseline & 256 & .004285 & 3.4022 & 3.4075 & 3.4081 \\
$\mu$P & 128 & .003073 & 3.6162 & 3.6162 & 3.6162 \\
$\mu$P & 256 & .002033 & 3.4017 & 3.4115 & 3.4178 \\
Both & 512 & .002378 & 3.2240 & 3.2291 & 3.2291 \\
\bottomrule
\end{tabular}
\end{center}
Figure~\ref{fig:p42-source-fits} shows the measured curves and the local
interpolations. Baseline's fitted optimum changes by $3.48\times$ between
widths 128 and 512, whereas $\mu$P's largest/smallest fitted optima differ
by $1.51\times$. However, better transfer does not imply uniformly better
tuned loss: $\mu$P wins the best sampled comparison by $.0113$ at width
128, while baseline wins by $.0041$ at width 256. The interpolated
width-256 minima are almost equal and reverse that tiny ordering.
There are no seed repeats, so such small differences are not evidence
of a statistically reliable winner.

\begin{figure}[htbp]
\centering
\includegraphics[width=\linewidth]{figures/p42_source_fits.png}
\caption{Source-width LR curves (circles), local log-LR quadratics and
fitted optima (stars). Bottom: fitted LR, fitted minimum loss, and
\emph{measured} loss at the transferred LR. Width increases from purple
to yellow; solid/dashed comparison lines distinguish baseline/$\mu$P.}
\label{fig:p42-source-fits}
\end{figure}

This is not a controlled change of training length alone: compared with
4.1's width test, these runs also have eight layers instead of two,
BF16 branch computations instead of fully FP32 computation, the course's
small shared initialization, and warmup, decay, WD and clipping.
It also uses a globally permuted training prefix, rather than the stress
test's contiguous prefix.
Any changed transfer behavior cannot be attributed only to more updates.

\paragraph{Takeaway.} $\mu$P improves width transfer, but the two policies
achieve similar losses after tuning; better transfer is not the same as
a uniformly better model.

\subsection{(c) Predict width 1,024}

For each prescription, the source-only optimal LRs at widths 128, 256
and 512 define a fitted rule
$\eta^*(n)=\eta_{512}(n/512)^p$. Numerical predictions and their source
measurements are saved before submitting or inspecting width-1,024
results. The predictions were frozen at 08:16:49 UTC on October 6:
\[
 \eta^*_{\rm baseline}(n)=.002350(n/512)^{-.9003},\qquad
 \eta^*_{\mu\mathrm P}(n)=.002162(n/512)^{-.1852}.
\]
Here the coefficient at 512 is estimated by the three-point log--log
regression, not forced to equal that width's individual fitted optimum.
They predict $.001259$ and $.001902$ at width 1,024, respectively.
Doubling width multiplies the fitted base LR by $.536$ for baseline
and $.880$ for $\mu$P. At the target width I test six rates per policy:
half the prediction, the prediction, twice the prediction, and
$\{.0015,.003,.006\}$. This includes direct $.003$ transfer and a local
bracketing sweep; loss gaps are relative to the best \emph{sampled}
target LR, not the fitted minimum.

All 12 target runs completed. Figure~\ref{fig:p42-target-fits} shows
the measured curves and bracketed local optima;
Figure~\ref{fig:p42-target-scaling} compares the frozen source laws with
the target results. Both gap columns below subtract that policy's best
sampled target loss.
\begin{center}
\small
\setlength{\tabcolsep}{4pt}
\begin{tabular}{lrrrrr}
\toprule
Policy & Predicted LR & Fitted LR & Best loss & Prediction gap & Direct gap \\
\midrule
Baseline & .001259 & .001018 & 3.110527 & .000000 & .118393 \\
$\mu$P & .001902 & .002425 & 3.101986 & .002921 & .000000 \\
\bottomrule
\end{tabular}
\end{center}

For the baseline, the prediction is the best of the six sampled rates;
direct $.003$ gives $3.228920$, so fitting width dependence improves loss
by $.118393$. For $\mu$P, direct $.003$ is already the best sampled rate;
the prediction gives $3.104907$, slightly worse by $.002921$. Fitting the
width trend therefore helps the baseline but does not improve $\mu$P
direct transfer in this experiment. The small $\mu$P difference is a
single-seed observation, not evidence of a universal disadvantage.

The predicted LRs are $23.7\%$ above and $21.6\%$ below their respective
local fitted optima. Those optima are interpolations, not additional
training runs: a prediction can be best sampled without exactly matching
the fitted optimum. Fitting all six points rather than only the local
three gives optima $.000987$ and $.002151$. This changes precise LR
estimates, especially for $\mu$P, but not the measured recipe comparisons.

\paragraph{Paired-seed check.} I repeated the two fixed $\mu$P recipes
at initialization seeds 43 and 44, keeping data order at seed 42 and
\emph{not} retuning LR. Prediction-minus-direct loss is $.002921$,
$.004650$ and $-.001099$ at seeds 42, 43 and 44. The mean is $.002157$,
with sample SD $.002950$. The small ranking therefore changes sign:
these repeats support similar performance, not a reliable advantage for
either $\mu$P recipe. Three seeds are still a small robustness check.

\begin{figure}[htbp]
\centering
\includegraphics[width=\linewidth]{figures/p42_target_fits.png}
\caption{Held-out width-1,024 loss--LR measurements and local quadratic
fits. Stars are interpolated optima, purple outlined circles identify
direct transfer, and teal hollow diamonds identify frozen predictions.
All target measurements use yellow, the largest width's color.}
\label{fig:p42-target-fits}
\end{figure}

\begin{figure}[htbp]
\centering
\includegraphics[width=\linewidth]{figures/p42_target_scaling.png}
\caption{Left: source fitted optima (stars), source-only power laws
(solid through width 512, dashed beyond it), frozen width-1,024 predictions
(teal hollow diamonds), and target fitted optima (yellow stars).
Right: measured prediction/direct-transfer loss gaps relative to each
policy's best sampled target LR. Marker colors match the preceding figure.}
\label{fig:p42-target-scaling}
\end{figure}

\paragraph{Takeaway.} Width fitting rescues baseline transfer. For
$\mu$P, direct transfer already works well; the fitted width rule does
not give a consistent gain across three initialization seeds.

\subsection{(d) Transfer across depth}

I fix width 512 and compare depths 4, 8 and 16 using the three 4.1 depth
policies, now with $r=L/8$. Depth eight is the reference: all policies
coincide and reuse the supplied curve. At depths four and sixteen, direct
$.003$ transfer and tuned loss are compared separately. I initially sweep
$\eta\in\{.0015,.003,.006\}$, then add $.00075$ wherever the best
sampled rate is at the lower boundary (21 new runs in total). Each
configuration's minimum is now bracketed; the fitted curves use the same
local log-LR quadratic as above.
At each fixed depth, all three policies start with identical stored
weights; I verified this through the shared trainer's actual model-builder
path on CPU. The residual and optimizer multipliers differ, not the
random initialization draws.

Figure~\ref{fig:p42-depth-fits} shows the curves, fitted optima and losses
at the directly transferred rate. The direct gap below is measured loss
at $.003$ minus the best measured loss for that same configuration.
It is not a difference from the interpolated minimum.
\begin{center}
\begin{tabular}{lrrrr}
\toprule
Policy & Depth & Fitted LR & Best measured loss & Direct gap \\
\midrule
All (reference) & 8 & .002378 & 3.2291 & .0000 \\
Standard $\mu$P & 4 & .002469 & 3.3156 & .0000 \\
Depth-$\mu$P & 4 & .001793 & 3.3258 & .0222 \\
CompleteP & 4 & .002517 & 3.3158 & .0000 \\
Standard $\mu$P & 16 & .001995 & 3.1739 & .0093 \\
Depth-$\mu$P & 16 & .002778 & 3.1686 & .0000 \\
CompleteP & 16 & .002057 & 3.1714 & .0046 \\
\bottomrule
\end{tabular}
\end{center}

At depth 16, Depth-$\mu$P improves \emph{both} comparisons: $.003$ is
already its best sampled rate, and its best measured loss is $.0053$
below standard $\mu$P's retuned loss. CompleteP reduces the direct gap
and improves the best measured loss by $.0025$. However, at depth four,
Depth-$\mu$P worsens direct transfer and has a higher tuned loss;
CompleteP and standard $\mu$P are nearly tied. Neither correction wins
universally. Across depths, the fitted LR spans $1.24\times$ for standard
$\mu$P, $1.55\times$ for Depth-$\mu$P and $1.22\times$ for CompleteP.

At depth 16 the fitted minimum ranks CompleteP ahead of Depth-$\mu$P
($3.1607$ versus $3.1681$), whereas the best \emph{measured} loss ranks
Depth-$\mu$P ahead. This small reversal depends on interpolation between
only three nearby measurements, so I do not treat interpolated minima
as a reliable ranking.

\paragraph{Paired-seed check at depth 16.} I repeated the seed-42 best
\emph{sampled} recipes at initialization seeds 43 and 44, without
retuning LR: standard $\mu$P and CompleteP use $.0015$, Depth-$\mu$P
uses $.003$, and data order remains seed 42. Negative differences below
mean improvement over standard $\mu$P at the same seed.
\begin{center}
\begin{tabular}{lrr}
\toprule
Seed & Depth-$\mu$P minus standard & CompleteP minus standard \\
\midrule
42 & $-.005296$ & $-.002515$ \\
43 & $-.011058$ & $-.001322$ \\
44 & $-.011635$ & $+.000305$ \\
\bottomrule
\end{tabular}
\end{center}
Depth-$\mu$P improves in all three seeds, with mean difference $-.009330$
and sample SD $.003505$. CompleteP's mean difference is only $-.001177$
(SD $.001415$), and its ranking changes sign. These are fixed-recipe
robustness checks, not seed-specific optimum sweeps or universal guarantees.

\begin{figure}[htbp]
\centering
\includegraphics[width=\linewidth]{figures/p42_depth_fits.png}
\caption{Longer-training depth comparison. Top: measurements (circles),
local quadratic fits and fitted optima (stars). Bottom: fitted optimal
base LR, fitted minimum loss and measured loss at directly transferred
LR $.003$. Colors order depth; depth 16 is yellow. Depth-eight measurements
are shared by all prescriptions.}
\label{fig:p42-depth-fits}
\end{figure}

\paragraph{Takeaway.} Depth-$\mu$P helps transfer and tuned loss at depth
16, with its selected recipe improving all three tested seeds, but it
does not help at depth four. CompleteP's small depth-16 advantage is
less consistent. Depth corrections are not automatic wins at every depth.

\subsection{(e) Diagnose the training regime}

At direct-transfer and best sampled rates, I collect logit RMS, normalized
feature movement and readout--movement alignment during the first five
updates and at the end. I also record pre-clipping gradient norm,
the embedding's share of squared gradient norm and the clipping
coefficient. In the baseline, the embedding scale $1/n$ can approach
$\sqrt{10^{-5}}$; then norm epsilon changes feature scale rather than
acting as a negligible numerical safeguard. Fixed WD also need not be
equally effective after changing LR, as Problem 2 demonstrated. A shifted
optimal LR alone therefore does not establish unstable $\mu$P behavior.

\paragraph{Width diagnostics, including the held-out model.}
Figures~\ref{fig:p42-width-features} and \ref{fig:p42-width-hidden}
track all four widths, at direct and best sampled rates. At width 1,024,
the logit, movement and alignment columns give update-five/end pairs; the raw residual
column gives only the end value. The baseline's best sampled rate is the
frozen prediction $.001259$, whereas $\mu$P's is direct $.003$.
\begin{center}
\small
\setlength{\tabcolsep}{4pt}
\begin{tabular}{lrrrr}
\toprule
Recipe & Logit RMS & Movement $M$ & $\omega_{\rm move}$ & Final residual RMS \\
\midrule
Baseline, direct & $1.001/4.077$ & $1.325/1.150$ & $.5000/.5011$ & 97.89 \\
Baseline, best & $.960/3.976$ & $1.328/1.357$ & $.4992/.5026$ & 16.75 \\
$\mu$P, direct=best & $.733/4.089$ & $1.315/1.394$ & $.5012/.5031$ & 21.26 \\
\bottomrule
\end{tabular}
\end{center}

Normalized features move substantially under both recipes; their movement
does not disappear at the held-out width. Yet direct baseline transfer
has a much larger raw residual than its retuned run, even though their
early normalized movements are almost identical. Looking only at the
normalized feature would miss that difference. $\mu$P's direct recipe
has controlled finite magnitudes and nonzero movement here, without an
LR change. This is finite-width evidence, not a proof of an infinite-width
stable regime.

At width 1,024, median update alignment at update five/end is
$.931/.673$ for direct baseline, $.918/.638$ for retuned baseline and
$.921/.641$ for direct $\mu$P. Along with the source-width heatmaps,
this supports strong early, weaker late and matrix-dependent alignment,
not one universal full-alignment assumption. The readout ratios remain
near $.5$ while features keep moving. They resemble the random reference
scale, but the finite-geometry caveat still prevents inferring independence.

\begin{figure}[htbp]
\centering
\includegraphics[width=\linewidth]{figures/p42_width_features.png}
\caption{Longer-training width diagnostics at direct and best sampled
rates, including held-out width 1,024 (yellow). Solid: direct $.003$;
dashed: best sampled, shown only when different. Columns show logit RMS,
final normalized-feature movement and initial-readout/movement alignment.
The horizontal axis labels actual updates and separates the first five
from later training; line segments join recorded checkpoints.}
\label{fig:p42-width-features}
\end{figure}

\begin{figure}[htbp]
\centering
\includegraphics[width=\linewidth]{figures/p42_width_features_hidden.png}
\caption{Width comparison of internal normalized movement and raw final
residual RMS, using the same rates, probes, colors and line styles as
Figure~\ref{fig:p42-width-features}. Baseline's larger raw width-1,024
stream at direct transfer is much smaller after retuning, despite similar
early normalized movements. Raw magnitude uses a log scale.}
\label{fig:p42-width-hidden}
\end{figure}

\paragraph{Depth diagnostics.}
Figures~\ref{fig:p42-depth-features} and \ref{fig:p42-depth-hidden}
compare the directly transferred and best sampled rates. At depth 16
and the best sampled rate, final-feature movement at update five/end
is $1.315/1.527$ for standard $\mu$P, $1.300/1.357$ for Depth-$\mu$P
and $1.315/1.525$ for CompleteP. These features change substantially;
none of these finite models is simply frozen at initialization.
First-block movement is smaller at update five ($.090$--$.171$),
but reaches $1.05$--$1.09$ by the end. Movement therefore depends on
both the layer and the training stage.

At the best sampled rates, median matrix-update alignment at depth 16
is approximately $.905$--$.911$ at update five and $.616$--$.621$ at
the final update. Early updates are much closer to the full-alignment
reference ($1$); late updates are weaker and closer to the random
reference ($.5$), but not equal to it. The feedforward down-projection
is especially weakly aligned at the end (median $.535$--$.537$).
The other depths show the same early-stronger/late-weaker pattern.
Meanwhile $\omega_{\rm move}$ remains near $.5$ at all depths despite
nonzero movement. As discussed in 4.0, this pooled readout statistic
does not establish independence of the initial readout and feature change.
No single exact full/no-alignment combination describes every layer
and training stage.

Raw scale also matters. At direct LR $.003$, the final residual RMS at
depths $4/8/16$ is $14.47/25.44/38.68$ under standard $\mu$P,
$34.49/25.44/16.00$ under Depth-$\mu$P, and
$29.16/25.44/18.85$ under CompleteP. Depth corrections change these
raw scales even when normalized movements look similar. These modest
depths do not establish asymptotic boundedness; conversely, an altered
raw scale or fitted optimum alone is not proof of instability.

\begin{figure}[htbp]
\centering
\includegraphics[width=.95\linewidth]{figures/p42_depth_features.png}
\caption{Longer-training depth diagnostics: logit RMS, final normalized
feature movement and $\omega_{\rm move}$ on fixed probes. Solid lines:
direct LR $.003$; dashed lines: best sampled LR, shown only when different.
Colors order depth, with depth 16 yellow. The horizontal axis separates
the first five updates from later training; connecting segments are not
extra measurements. Undefined zero-movement alignment is omitted.}
\label{fig:p42-depth-features}
\end{figure}

\begin{figure}[htbp]
\centering
\includegraphics[width=\linewidth]{figures/p42_depth_features_hidden.png}
\caption{Internal normalized-feature movement at the first, midpoint and
last blocks, alongside raw final residual RMS (log scale). Each movement
is measured from that same layer's initial feature on the same fixed
inputs. Colors and direct/best line styles match
Figure~\ref{fig:p42-depth-features}. Normalized movement and raw magnitude
are separate diagnostics.}
\label{fig:p42-depth-hidden}
\end{figure}

\paragraph{Normalization, clipping and weight decay.}
More concretely, RMSNorm divides by
$\sqrt{\operatorname{mean}(x_i^2)+\epsilon}$, not just the feature RMS.
If the embedding RMS is $s$, the first normalized-feature RMS is
approximately $s/\sqrt{s^2+10^{-5}}$. The baseline's $s\approx1/n$
therefore becomes increasingly attenuated with width; $\mu$P's
$s\approx1/512$ keeps this particular effect approximately fixed.
Later normalized layers can still be near unit RMS, hiding this
first-layer difference.

Figure~\ref{fig:p42-initial-norm} makes this difference visible before
any training: the baseline's first normalized-feature RMS is $.925$,
$.773$, $.520$ and $.291$ at widths 128, 256, 512 and 1,024;
$\mu$P stays approximately $.520$ at all four.
This is a measured width-dependent normalization
effect, not just a hypothetical epsilon concern.
At the source LR and width 128, $\mu$P also has a larger logit RMS at
update five ($1.843$ versus $.970$), yet a lower final loss ($3.6162$
versus $3.6940$). A larger early transient alone does not identify a worse
final training recipe.

\begin{figure}[htbp]
\centering
\includegraphics[width=.85\linewidth]{figures/p42_initial_normalization.png}
\caption{Measured first-RMSNorm output at initialization. Baseline's
embedding scale shrinks with width while norm epsilon stays fixed;
$\mu$P's fixed-reference embedding scale holds this output nearly constant.
These are initialization measurements, independent of the tuned LR.}
\label{fig:p42-initial-norm}
\end{figure}

There is also an explicit regularization coupling: a hidden $\mu$P weight
is shrunk by $1-(\eta_t/m)\lambda$ per update, compared with
$1-\eta_t\lambda$ in the baseline. Depth-$\mu$P adds another $r^{-1/2}$
to the block rate. Thus fixed $\lambda=.1$ does not hold hidden-weight
shrinkage fixed across prescriptions. Embeddings and norms have no WD.
These are competing explanations to check against the diagnostics, not
causal conclusions from a shifted optimum alone.
At width 1,024 and direct LR, the initial gradient norm is $33.84$ in
baseline and $11.84$ in $\mu$P; the embedding accounts for $99.0\%$ and
$96.4\%$ of their squared norms. Clipping multiplies these gradients by
$.0296$ and $.0845$, respectively. By the final update both have clipping
coefficient one and embedding shares below $3\%$.
Figure~\ref{fig:p42-width-clipping} shows that this strongly
embedding-dominated, clipped early stage is not the whole training regime.
The single-norm intervention in Problem~\ref{sec:p5-norm} directly tests
this normalization effect: removing the first layer's attenuation changes
initial feature RMS substantially, but does not consistently improve loss.

\begin{figure}[htbp]
\centering
\includegraphics[width=\linewidth]{figures/p42_width_features_clipping.png}
\caption{Gradient and clipping diagnostics for the same width recipes
as Figure~\ref{fig:p42-width-features}. Left: full pre-clipping norm;
middle: embedding's fraction of squared norm; right: applied clipping
coefficient. Early gradients are predominantly embedding gradients and
strongly clipped; clipping is absent at the final checkpoint. These are
gradient diagnostics, not direct measurements of Adam's parameter step.}
\label{fig:p42-width-clipping}
\end{figure}

Unlike the independent stress-test draws, the course's shared Gaussian
embedding/readout initialization also means that the initial readout and
initial features are not strictly independent. The theory's independent-
initialization reference is therefore not literally this longer-training
setup; the measured alignment and movement must be interpreted directly.
For Adam, a smaller clipping coefficient also does not imply the same
proportional reduction in the parameter update: uniformly scaling a
gradient largely cancels between its first moment and square-root second
moment when the stabilizer is negligible. Changing clipping across updates
can still change those running moments. I therefore measure actual
parameter updates rather than infer their sizes from gradient norm alone.
Finally, the early checkpoints occur during warmup and the final update
occurs near the end of linear LR decay. The stage comparison therefore
includes a changing step size and a changing model, not just more training
at one fixed LR. The alignment measurements describe this recipe's actual
updates; they do not isolate which factor caused the change.

\paragraph{Takeaway.} Alignment changes by matrix and training stage,
while features keep moving. Judge transfer with losses \emph{and} raw
signal diagnostics: neither a shifted LR nor a larger early transient
alone proves instability. Fixed WD, norm epsilon and clipping can change
the comparison without invalidating the entire $\mu$P recipe.

\subsection{(f) Optional: Muon and width transfer}

Muon treats a whole hidden weight matrix together, rather than dividing
each gradient coordinate by Adam's running RMS. I reuse the course's
implementation unchanged: Nesterov momentum $.95$, five BF16
Newton--Schulz iterations to approximately orthogonalize the momentum
direction, and the supplied rectangular-matrix correction
$\sqrt{\max(1,d_{\rm out}/d_{\rm in})}$. Embedding, readout and
normalization parameters still use auxiliary AdamW with $(.9,.95)$.
As in the course implementation, WD applies to hidden matrices and the
readout, but not to the embedding or norm weights.
The implementation is adapted from
\href{https://github.com/KellerJordan/Muon/tree/f98f1cacc0263b04290753e32be8d498c1efc806}{Keller Jordan's Muon code}.

An important difference from the required AdamW experiment is that
\emph{hidden Muon LR stays at the base value across widths}; I do not
copy Adam's $1/m$ multiplier. An ideal orthogonalized update $Q$ has
operator norm one, so
\[
\operatorname{RMS}(Qx)\leq
\sqrt{d_{\rm in}/d_{\rm out}}\operatorname{RMS}(x).
\]
With the matrix aspect ratio fixed, this bound does not grow with width.
The supplied rectangular correction is also width-independent at fixed
aspect ratio. This motivates a constant hidden base LR, rather than an
extra inverse-width factor. This optimizer-specific distinction is
also discussed by
\href{https://proceedings.neurips.cc/paper_files/paper/2025/file/bdcd9f6327db5877dee502cdec183159-Paper-Conference.pdf}{Qiu et al. (2025), Table 1 and Appendix G}.
The real Newton--Schulz update is approximate; the bound explains the
scaling choice, not an exact guarantee for every measured update.

The reference pilot's best pair is hidden LR $.02$ and auxiliary AdamW
LR $.0015$, with loss $3.107776$. Holding hidden LR $.02$ fixed, the
smaller auxiliary LR $.0003$ gives $3.157971$: the auxiliary choice matters.

I compare the same baseline and $\mu$P initialization/readout policies
as above, now with Muon in the hidden matrices. The auxiliary AdamW LR
is separately tuned at reference width 512: hidden LR
$\{.01,.02,.04\}$ crossed with auxiliary LR $\{.0003,.0015\}$.
After choosing the best reference pair, auxiliary LR $.0015$ is held fixed
for source-width sweeps and the held-out width. Only two auxiliary rates
were tested, so this is a conditional choice, not a fully optimized auxiliary LR.
All runs keep 153.6M
tokens, batch 64, WD $.1$, linear decay, 1\% warmup and clipping at one.
The two policies coincide at width 512, so those reference runs are shared.
The initial hidden-LR grid at each source width is
$\{.005,.01,.02,.04,.08\}$; additional endpoint points are added only
if needed to bracket a minimum. As in the required part, the loss fit is a
quadratic in log LR through the best measured point and its two neighbors.
The direct-transfer LR is the best \emph{Muon} reference LR, not AdamW's
$.003$: changing optimizer does not imply that its best base LR is unchanged.
All source sweeps completed, giving the optima in
Figure~\ref{fig:p42f-source}:
\begin{center}
\begin{tabular}{lrr}
\toprule
Width & Baseline fitted LR & $\mu$P fitted LR \\
\midrule
128 & .035666 & .032364 \\
256 & .030139 & .028663 \\
512 & .023222 & .023222 \\
\bottomrule
\end{tabular}
\end{center}
The source-only power laws are
\[
\eta^*_{\mathrm{base}}(n)=.023582(n/512)^{-.3095},\qquad
\eta^*_{\mu\mathrm{P}}(n)=.023569(n/512)^{-.2395}.
\]
Their prefactors are least-squares estimates, not forced to equal the
width-512 quadratic optimum. Predictions at width 1,024 are $.019028$
and $.019964$, respectively. These were frozen before any held-out
submission. Direct transfer uses $.02$; the $\mu$P prediction differs
from it by only $.18\%$, so these are nearly the same recipe.
As a fitting-sensitivity check, using all five source points rather than
local neighbors gives about $.0208$ for both held-out LRs. These are
not separately tested predictions and do not replace the frozen values.

\begin{figure}[htbp]
\centering
\includegraphics[width=\linewidth]{figures/p42f_muon_source.png}
\caption{Muon source loss--LR curves, local quadratic optima (stars),
and source-only width laws. Solid lines span source widths; dashed lines
extrapolate to the held-out width, with frozen predictions as hollow diamonds.
Yellow stars mark held-out fitted optima. Dotted vertical lines mark direct LR $.02$.}
\label{fig:p42f-source}
\end{figure}

The largest-to-smallest source optimum ratio is $1.536$ for the baseline
and $1.394$ for $\mu$P, versus $3.484$ and $1.512$ with AdamW.
Muon therefore makes the baseline substantially less width-sensitive;
$\mu$P gives a smaller additional improvement. At direct LR $.02$, the
loss gaps to the best sampled source point are $.008763$ and $.002842$
for baseline widths 128 and 256, versus $.003172$ and $.000679$ for
$\mu$P. Figure~\ref{fig:p42f-performance} also shows lower tuned losses
for $\mu$P at these two widths. These are conditional comparisons with
one fixed auxiliary LR and WD, not proof of a universally best prescription.
At the shared reference width 512, Muon's best sampled loss is $3.107776$
versus AdamW's $3.229134$: a $.121358$ improvement at the same token budget.
This is not a compute-matched or equally searched optimizer comparison.
Cross-optimizer comparisons also retain the LR--WD coupling from Problem 2:
different hidden LRs at fixed WD give different multiplicative weight
shrinkage. These runs do not isolate orthogonalization from that difference.

\begin{figure}[htbp]
\centering
\includegraphics[width=\linewidth]{figures/p42f_muon_source_performance.png}
\caption{Muon tuned performance (left, fitted stars) and measured loss
at the transferred reference LR $.02$ (right, circles). Source width 512
is shared by both prescriptions and is yellow.}
\label{fig:p42f-performance}
\end{figure}

Figures~\ref{fig:p42f-align-base} and \ref{fig:p42f-align-mup} show every
matrix's update alignment at direct LR $.02$, including held-out width
1,024. Across source widths, the median is near $.51$ at updates 2 and 5
and $.50$ at the final update; individual
readout/rectangular-matrix values can be larger (up to about $.91$ at
those checkpoints and $.97$ elsewhere in training).
This differs from assuming maximally aligned hidden updates. But $.5$
need not mean statistical independence: for a square orthogonal update
$Q$, $\operatorname{RMS}(QX)=\operatorname{RMS}(X)$ and
$\operatorname{RMS}(Q)=1/\sqrt n$, hence
\[
\frac{\operatorname{RMS}(QX)}{\operatorname{RMS}(Q)\operatorname{RMS}(X)}
=\sqrt n,\qquad \alpha_{\mathrm{upd}}=.5
\]
for \emph{any} $X$. Real updates are only approximately orthogonal and
include WD; this identity explains the measurement without establishing
an exact theorem for every matrix shape.

\begin{figure}[htbp]
\centering
\includegraphics[width=\linewidth]{figures/p42f_muon_alignment_baseline.png}
\caption{Muon baseline: update alignment for the three source widths,
held-out width 1,024, and
all attention, feedforward and readout matrices. Checkpoint rows are
ordinal, not evenly spaced in training time; white denotes zero-update
undefined ratios, including the first warmup update.}
\label{fig:p42f-align-base}
\end{figure}
\begin{figure}[htbp]
\centering
\includegraphics[width=\linewidth]{figures/p42f_muon_alignment_mup.png}
\caption{Muon with $\mu$P: the same per-matrix measurements, direct LR
and color scale as Figure~\ref{fig:p42f-align-base}.}
\label{fig:p42f-align-mup}
\end{figure}

Figure~\ref{fig:p42f-diagnostics} compares feature movement. At the three
source widths, final normalized movement is $1.52$--$2.00$ for the baseline
and $1.40$--$1.52$ for $\mu$P. Readout--feature-change alignment stays
near $.50$; again, this is a measured ratio, not proof of independence.
Raw residual RMS peaks at $9.80$--$10.16$ around update 234 before
falling to $2.88$--$3.94$ at the end. These finite transients do not by
themselves establish asymptotic stability or instability.
Clipping still affects the source comparison: the initial embedding contributes
$90$--$99\%$ of squared gradient norm, and clipping coefficients range
from $.034$ to $.263$. At the end, coefficients are one at all sources.
Thus Muon's update normalization does not remove the embedding, norm
epsilon or auxiliary-optimizer confounders.
At held-out width 1,024, final normalized movement is $1.330$ for the
baseline and $1.594$ for $\mu$P; hidden update-alignment medians are
$.508$ and $.506$, and readout movement alignment is about $.501$.
Raw residual RMS peaks at $10.60$ and $10.34$, respectively. Initial
clipping coefficients are $.030$ and $.084$, becoming one at the end.
Both larger models still change features substantially; these finite
diagnostics do not prove an infinite-width training regime.

\begin{figure}[htbp]
\centering
\includegraphics[width=\linewidth]{figures/p42f_muon_diagnostics.png}
\caption{Muon source and held-out diagnostics at fixed hidden LR $.02$ and auxiliary
LR $.0015$: normalized feature movement, readout movement alignment and
raw residual RMS. The readout-alignment axis is zoomed around $.5$;
its dotted line marks $.5$. Undefined zero-movement points are omitted.
Yellow denotes held-out width 1,024, the largest width.}
\label{fig:p42f-diagnostics}
\end{figure}

\paragraph{Held-out width 1,024.}
All twelve target runs completed: six rates per policy, consisting of
the frozen prediction, direct reference LR $.02$, and half/double
neighbors of each. Figure~\ref{fig:p42f-target} shows the measured
curves and target fits. Because predicted/direct LRs and their neighbors
are close, I use a least-squares quadratic in $\log_2(\eta/.003)$ over
\emph{all six target points}, rather than a local interpolant dominated
by nearly coincident LRs. This fitting choice was recorded before the
complete target curves were available. Loss gaps below use actual measured
losses, not fitted minima:
\begin{center}
\small
\begin{tabular}{lrrrr}
\toprule
Policy & Fitted LR & Best sampled loss & Prediction gap & Direct gap \\
\midrule
Baseline & .019503 & 2.974521 & .000000 & .000597 \\
$\mu$P & .020113 & 2.976238 & .000000 & .001460 \\
\bottomrule
\end{tabular}
\end{center}
Both predictions are the best sampled rates. Direct transfer is already
close to the minimum: its losses are $2.975117$ and $2.977699$.
Unlike AdamW's baseline, where direct transfer incurred a $.118393$ gap,
Muon needs little correction at this held-out width. The $\mu$P prediction
and control differ by only $.18\%$ in LR, so their small single-seed
loss difference does not establish a broadly useful extra scaling correction.

The precise fitted optimum depends on the fitting method. Local
three-point interpolation instead gives $.017585$ and $.014284$;
especially for $\mu$P, tiny horizontal gaps turn small loss differences
into a sharply curved interpolant. The full-grid fit is a descriptive
estimate, not a confidence interval or another training run. The measured
recipe comparisons do not depend on which curve is fitted.

\begin{figure}[htbp]
\centering
\includegraphics[width=\linewidth]{figures/p42f_muon_target.png}
\caption{Muon at held-out width 1,024. Left/middle: measured losses
(yellow circles), full-grid least-squares quadratics and fitted optima
(yellow stars). Purple outlined circles mark direct LR $.02$; teal hollow
diamonds mark the frozen predictions. Right: their measured loss gaps
relative to each policy's best sampled target point, using the same colors.}
\label{fig:p42f-target}
\end{figure}

\paragraph{Paired-seed check.} At fixed source-selected hidden LR $.02$
and auxiliary LR $.0015$, I repeated both policies at initialization seeds
43 and 44, keeping data order at 42 and reusing the seed-42 controls.
These compare fixed recipes, not independently retuned optima:
\begin{center}
\begin{tabular}{lrrr}
\toprule
Seed & Baseline loss & $\mu$P loss & $\mu$P minus baseline \\
\midrule
42 & 2.975117 & 2.977699 & $+.002581$ \\
43 & 2.974576 & 2.976628 & $+.002053$ \\
44 & 2.974340 & 2.975357 & $+.001017$ \\
\bottomrule
\end{tabular}
\end{center}
The mean difference is $+.001884$, with sample SD $.000796$.
The small source advantage for $\mu$P does not persist at the held-out
fixed recipe: the baseline is slightly better at all three tested seeds.
This is not a universal policy ranking; auxiliary LR and WD were not
retuned, data order is fixed, and three initialization seeds remain a
small check.

\paragraph{Takeaway.} Muon transfers LR well; width fitting adds little.
Adding $\mu$P does not improve this held-out fixed recipe in three seeds.

\section{Problem 5: Optional prediction question}
\label{sec:p5-norm}

\paragraph{Question: when is a numerical stabilizer not negligible?}
Use the width-1,024 baseline from 4.2, with eight layers, the same
initialization seed and data, LR $.003$ and 153.6M tokens. Change
\emph{only} the first block's input RMSNorm epsilon from $10^{-5}$ to
$10^{-8}$; all other normalization epsilons stay unchanged.
Measure RMS on the shared logger's fixed eight validation sequences
(8,192 token positions); final validation loss uses all 1,000 sequences.
Which statements hold at initialization? Select all that apply.
\begin{enumerate}
\item[A.] First-normalized feature RMS lies between $.95$ and $1.05$.
\item[B.] Its RMS is between $3.1$ and $3.7$ times the original value.
\item[C.] All stored initialization weights remain identical at a fixed seed.
\item[D.] Initial logits remain numerically unchanged.
\end{enumerate}

\paragraph{Recorded prediction.} I predicted A, B and C before
submitting the intervention. For embedding RMS $s$, the initial
normalized RMS is approximately $s/\sqrt{s^2+\epsilon}$. Here
$s\approx1/1024$, so $\epsilon=10^{-5}$ strongly attenuates the feature;
$10^{-8}$ is much smaller than $s^2$, bringing normalized RMS close to
one. The code changes a non-parameter scalar, not an initialization
draw, but changes the first forward branch and hence generally the logits.

The executable intervention, recorded prediction and launch script are
in \texttt{experiments/a2/p5\_normalization.py} (62 lines, below the
100-line limit). The focused model change is:
\begin{verbatim}
class FirstNormLM(ParameterizedLM):
    def __init__(self, *args, first_norm_epsilon=1e-8, **kwargs):
        super().__init__(*args, **kwargs)
        self.model.layers[0].input_layernorm.variance_epsilon = (
            first_norm_epsilon)
\end{verbatim}
The shared trainer selects this subclass through its existing model-builder
interface; no optimizer or training-loop change is needed.
A CPU check at the actual width 1,024 and depth eight (seed 43) confirmed
exact equality of all 91 stored parameter tensors and exactly one changed
norm epsilon. I use paired initialization
seeds 42, 43 and 44, reusing the original seed-42 control, so five new
runs test the question and any downstream loss change. All five completed.
The measured first-normalized RMS is $.9946$ at every seed, versus
$.2908$--$.2916$ before the intervention: a $3.411$--$3.420\times$ increase.
Initial logit RMS also changes at every seed. Thus the empirical answer
key is \textbf{A, B and C}, matching the recorded prediction; D is false.
The direct formula predicts new RMS within about $10^{-5}$ of the
measured value. Stored-weight equality follows from the shared initializer
and the separately verified constructor control, not from similar losses.

\begin{center}
\begin{tabular}{lrrr}
\toprule
Seed & Original final loss & New final loss & New minus original \\
\midrule
42 & 3.228920 & 3.258916 & $+.029995$ \\
43 & 3.212946 & 3.231405 & $+.018459$ \\
44 & 3.231832 & 3.210068 & $-.021763$ \\
\bottomrule
\end{tabular}
\end{center}
Final loss does not improve consistently: mean difference $+.008897$
with sample SD $.027172$. The robust initialization-scale prediction is
different from a reliable performance improvement. LR and WD were not
retuned after changing normalization.

\paragraph{Conceptual takeaway.} Epsilon is harmless only when it is
small relative to the variance being normalized. Otherwise, changing it
changes signal scale even with identical stored weights---but removing
that attenuation is not automatically a better training recipe.

\end{document}

